\documentclass[11pt]{article}
\pdftrailerid{}
\usepackage[margin=1.05in]{geometry}
\usepackage{amsmath,amssymb,amsthm}
\ifdefined\pdfglyphtounicode
  \pdfgentounicode=1
  \pdfglyphtounicode{tfm:cmsy10/mapsto}{007C}
  \pdfglyphtounicode{tfm:cmex10/parenleftbig}{0028}
  \pdfglyphtounicode{tfm:cmex10/parenrightbig}{0029}
  \pdfglyphtounicode{tfm:cmex10/bracketleftbig}{005B}
  \pdfglyphtounicode{tfm:cmex10/bracketrightbig}{005D}
  \pdfglyphtounicode{tfm:cmex10/vextendsingle}{23D0}
  \pdfglyphtounicode{tfm:cmex10/parenleftbigg}{0028}
  \pdfglyphtounicode{tfm:cmex10/parenrightbigg}{0029}
  \pdfglyphtounicode{tfm:cmex10/bracketleftbigg}{005B}
  \pdfglyphtounicode{tfm:cmex10/bracketrightbigg}{005D}
  \pdfglyphtounicode{tfm:cmex10/floorleftbigg}{230A}
  \pdfglyphtounicode{tfm:cmex10/floorrightbigg}{230B}
  \pdfglyphtounicode{tfm:cmex10/braceleftbigg}{007B}
  \pdfglyphtounicode{tfm:cmex10/bracerightbigg}{007D}
  \pdfglyphtounicode{tfm:cmex10/parenleftBigg}{0028}
  \pdfglyphtounicode{tfm:cmex10/parenrightBigg}{0029}
  \pdfglyphtounicode{tfm:cmex10/bracketleftBigg}{005B}
  \pdfglyphtounicode{tfm:cmex10/bracketrightBigg}{005D}
  \pdfglyphtounicode{tfm:cmex10/braceleftBigg}{007B}
  \pdfglyphtounicode{tfm:cmex10/bracerightBigg}{007D}
  \pdfglyphtounicode{tfm:cmex10/bracketlefttp}{23A1}
  \pdfglyphtounicode{tfm:cmex10/bracketrighttp}{23A4}
  \pdfglyphtounicode{tfm:cmex10/bracketleftbt}{23A3}
  \pdfglyphtounicode{tfm:cmex10/bracketrightbt}{23A6}
  \pdfglyphtounicode{tfm:cmex10/bracelefttp}{23A7}
  \pdfglyphtounicode{tfm:cmex10/bracerighttp}{23AB}
  \pdfglyphtounicode{tfm:cmex10/braceleftbt}{23A9}
  \pdfglyphtounicode{tfm:cmex10/bracerightbt}{23AD}
  \pdfglyphtounicode{tfm:cmex10/braceleftmid}{23A8}
  \pdfglyphtounicode{tfm:cmex10/bracerightmid}{23AC}
  \pdfglyphtounicode{tfm:cmex10/braceex}{23AA}
  \pdfglyphtounicode{tfm:cmex10/circleplusdisplay}{2A01}
  \pdfglyphtounicode{tfm:cmex10/summationtext}{2211}
  \pdfglyphtounicode{tfm:cmex10/producttext}{220F}
  \pdfglyphtounicode{tfm:cmex10/integraltext}{222B}
  \pdfglyphtounicode{tfm:cmex10/summationdisplay}{2211}
  \pdfglyphtounicode{tfm:cmex10/productdisplay}{220F}
  \pdfglyphtounicode{tfm:cmex10/integraldisplay}{222B}
  \pdfglyphtounicode{tfm:cmex10/tildewide}{02DC}
  \pdfglyphtounicode{tfm:cmex10/radicalbig}{221A}
  \pdfglyphtounicode{tfm:cmex10/radicalbigg}{221A}
  \pdfglyphtounicode{tfm:cmex10/radicalBigg}{221A}
  \pdfglyphtounicode{tfm:cmex8/summationtext}{2211}
  \pdfglyphtounicode{tfm:cmex8/integraltext}{222B}
  \pdfglyphtounicode{tfm:cmex8/tildewide}{02DC}
\fi

\usepackage{microtype}
\usepackage{flafter}
\usepackage[colorlinks=true,linkcolor=blue,citecolor=blue,urlcolor=blue]{hyperref}
\hypersetup{
  pdftitle={A counterexample to integer-degree harmonic dimension comparison},
  pdfauthor={OpenAI},
  pdfsubject={Polynomial-growth harmonic functions on manifolds with nonnegative Ricci curvature},
  pdfkeywords={harmonic functions, polynomial growth, nonnegative Ricci curvature, dimension comparison}
}
\numberwithin{equation}{section}
\newtheorem{theorem}{Theorem}[section]
\newtheorem{proposition}[theorem]{Proposition}
\newtheorem{lemma}[theorem]{Lemma}
\newtheorem{corollary}[theorem]{Corollary}
\theoremstyle{remark}
\newtheorem{remark}[theorem]{Remark}
\newcommand{\R}{\mathbb R}
\newcommand{\C}{\mathbb C}
\newcommand{\N}{\mathbb N}
\newcommand{\Ric}{\operatorname{Ric}}
\newcommand{\tr}{\operatorname{tr}}
\newcommand{\Id}{I}
\newcommand{\Hk}{\mathcal H}
\newcommand{\tf}{\mathrm{tf}}
\newcommand{\vol}{\operatorname{vol}}
\newcommand{\SL}{\operatorname{SL}}
\newcommand{\SO}{\operatorname{SO}}
\newcommand{\sgn}{\operatorname{sgn}}
\newcommand{\diag}{\operatorname{diag}}
\title{A counterexample to integer-degree\\harmonic dimension comparison}
\author{OpenAI}
\date{September 25, 2026}

\begin{document}
\maketitle
\begin{abstract}
For some even $n\ge8$ and integer $k\ge2$, we construct a complete smooth
metric on $\mathbb R^n$ with nonnegative Ricci curvature whose space of real
harmonic functions of pointwise polynomial growth at most $k$ has dimension
larger than the Euclidean harmonic-polynomial dimension. The metric is
Euclidean near the origin, has asymptotic volume ratio strictly between zero
and one, and has nonunique tangent cones at infinity. This answers the integer-degree form of Yau's dimension comparison
question in the negative.
\end{abstract}

\section{Introduction}\label{sec:introduction}

For a complete connected smooth Riemannian manifold $(M^n,g)$ without
boundary, let $\mathcal H_d(M,g)$ be the real vector space of harmonic
functions satisfying
\[
 |u(x)|\leq C_u(1+d_g(o,x))^d\qquad(x\in M)
\]
for some finite constant $C_u$, where $d\geq0$ and $o\in M$ is a base
point. Changing $o$ does not change the space. Write
$h_d(M,g)=\dim_{\mathbb R}\mathcal H_d(M,g)$. On Euclidean space, an
entire harmonic function with integer growth at most $k$ is a harmonic
polynomial of degree at most $k$. Consequently
\begin{equation}\label{eq:euclidean-comparator}
 h_k(\mathbb R^n,g_{\mathrm E})
 =\binom{n+k-1}{k}+\binom{n+k-2}{k-1}\qquad(k\geq1).
\end{equation}
The integer-degree comparison problem asks whether $\operatorname{Ric}_g
\geq0$ forces $h_k(M,g)\leq h_k(\mathbb R^n,g_{\mathrm E})$ for every
integer $k\geq1$. The condition defining $\mathcal H_k$ is pointwise at
every distance. Bounds only along a sequence of spheres do not establish it.

\begin{theorem}\label{thm:main}
There are an even integer $n\geq8$, an integer $k\geq2$, and a complete
smooth metric $g$ on $\mathbb R^n$ such that $\operatorname{Ric}_g\geq0$
and
\[
 h_k(\mathbb R^n,g)>h_k(\mathbb R^n,g_{\mathrm E}).
\]
The metric is Euclidean near the origin, and its asymptotic volume ratio
satisfies
\[
 0<\lim_{R\to\infty}
 \frac{\operatorname{vol}_g B_g(0,R)}{\omega_nR^n}<1,
\]
where $\omega_n$ is the volume of the Euclidean unit ball.
\end{theorem}

The exact finite spectral selection in Appendix~\ref{sec:finite-selection}
permits the choices \(n=16\) and \(k=50000\).

Theorem~\ref{thm:main} answers the integer-degree form of Yau's dimension
comparison question in the negative. Yau's questions on polynomial-growth harmonic functions distinguish finite
dimensionality from the sharper Euclidean comparison~\cite{Yau,LiTam}.
Li and Tam proved the sharp linear-growth bound $h_1(M,g)\leq n+1$ under
nonnegative Ricci curvature~\cite{LiTam}. Colding and Minicozzi proved
finite dimensionality at every fixed degree and later obtained bounds of
the optimal order $d^{n-1}$ as $d$ grows~\cite{CM97,CM98}. Those estimates
leave open the exact Euclidean count at a particular integer degree.

Donnelly constructed counterexamples to the analogous comparison at a real
degree strictly between one and two in dimensions at least
five~\cite{Donnelly,CaiLai}. That fact alone does not decide the integer
endpoint two: the Euclidean space of quadratic harmonic polynomials is
larger than its space at a degree between one and two. Cai and Lai make
this distinction explicit and establish sharp comparison with the
additional hypothesis of local conformal flatness~\cite{CaiLai}.
The theorem above addresses the integer question under the Ricci
hypothesis alone.

Cone-spectral results explain another boundary of the problem. For complete
manifolds with nonnegative Ricci curvature, maximal volume growth, and a
unique tangent cone at infinity, Huang's
Theorem~1.6 relates the harmonic dimension to the spectrum of the cone
link~\cite{Huang}. Our metric instead has round and nonround subsequential
cone limits. Nonunique cones with positive Ricci curvature and Euclidean
volume growth already occur in Perelman's construction~\cite{Perelman};
Colding and Naber developed a method for slowly varying link
families with a common volume form~\cite{ColdingNaber}. We use the same geometric possibility of a
moving link, while proving the curvature estimates for our particular
deformations. Xu's nonunique-cone three-circles theorem requires a comparison
exponent outside the union of cone-degree spectra~\cite[Theorem~3.4]{Xu}.
A separate geometric conclusion of our construction is that this union
contains a neighborhood of the selected integer $k$, so the theorem cannot
be applied at $k$ or at sufficiently nearby exponents.

\subsection*{Why a moving link can change the count}

Put $n=m+1$. On the metric cone $dr^2+r^2h$ over an $m$-dimensional link,
an angular eigenfunction of the nonnegative Laplacian with eigenvalue
$\lambda$ has a growing homogeneous harmonic mode $r^d\phi$, where
\[
 d(d+m-1)=\lambda,\qquad d\geq0.
\]
For one fixed link, each angular direction therefore carries one fixed
growth rate. Our construction lets a harmonic function use several rates
in succession. The challenge is to arrange that succession exactly:
an arbitrarily small component in an unwanted faster direction can dominate
after a sufficiently long interval.

The links are volume-normalized Berger metrics on odd spheres
$S^m$, $m=2s-1$. Their distinguished Hopf direction is allowed to change.
Differentiation along this circle direction is the Hopf generator; on
complexified harmonics its eigenvalues are \(\mathrm i\) times the
integer Hopf charges. Every angular Laplacian nevertheless preserves the same finite-dimensional
space $V_l$ of round spherical harmonics of degree $l$; write
$N_l=\dim_{\mathbb R}V_l$. At a fixed nonround link, write its growing rates
on $V_l$ in increasing order as $d_{l,1},\ldots,d_{l,N_l}$. The spectral
calculation supplies a finite $L\geq2$ and integers $0\leq M_l\leq N_l$
for $2\leq l\leq L$ such that
\begin{equation}\label{eq:overview-selection}
 \sum_{l=2}^L M_l>h_k(\mathbb R^{m+1},g_{\mathrm E}),
 \qquad
 \frac1{M_l}\sum_{i=1}^{M_l}d_{l,i}<k\quad(M_l>0).
\end{equation}
Thus the count will be large enough if each selected harmonic function can
spend asymptotically equal amounts of logarithmic radius at all the selected
rates in its block. The inequality for the mean is strict; individual rates
may exceed $k$.

The source of \eqref{eq:overview-selection} is the distribution of Hopf
charges inside $V_l$. The squared charges divided by $l^2$ tend to a
beta distribution. The mean of the lowest fraction of rates is described by
an integral of its decreasing quantile. In a sufficiently large odd link
dimension, a weighted quantile inequality shows that selecting directions by
their mean produces a strict surplus over the Euclidean count. This fixes the
link parameters, the integer $k$, and a finite degree range before the
radial construction begins. An exact-integer computation gives a separate
finite selection at specific parameters; its role is to verify that finite
spectral input, rather than the subsequent geometric or transmission proof.

\subsection*{From angular control to entire harmonic functions}

Short deformations near a round link provide the exchanges of directions.
For every fixed finite range $2\leq l\leq L$, the trace-free squares of
the Hopf generators, using all orthogonal complex structures, generate
\[
 \bigoplus_{l=2}^L\mathfrak{sl}(V_l).
\]
The real form and the direct sum matter: one must choose the same finite
sequence of geometric pulses while independently arranging a prescribed
signed permutation in each $V_l$. A finite product of exponentials realizes
these permutations at a parameter value where the word map has invertible
differential. The transmission argument below shows that this algebraic
freedom survives the harmonic equation.

Write the metric as $g=dr^2+r^2\gamma(t)$ with $t=\log r$. A harmonic
function in $V_l$ is governed by a second-order matrix equation. Smoothness
at the Euclidean center selects a distinguished regular solution space.
Its logarithmic derivative $P_l$ satisfies a matrix Riccati equation;
positive barriers keep $P_l$ bounded and symmetric. A long round rest makes
$P_l$ approach a scalar matrix, independently of all earlier admissible
pulses. During the next control interval, which is short relative to the
following dwell, the determinant-normalized value transfer converges in
$C^1$ to the finite algebraic word as the period index tends to infinity.

The pulse-end approximation is only an intermediate fact: the outgoing
radial derivative can still inject a small component into a faster direction
during the following long leg, on which the angular eigenvectors remain
fixed. We factor out the positive diagonal transfer \(D_{l,j}\) obtained by entering that entire leg with the round
growing derivative. The remaining multiplier of the derivative error is
bounded independently of the leg's length and anisotropy
(Lemma~\ref{lem:leg}). Thus the full-period map, with this diagonal
factor removed and its determinant normalized, is \(C^1\)-close to the ideal
word, uniformly over all preceding controls. Its local inverse tunes every complete period exactly to
\(s_{l,j}D_{l,j}\Pi_l\), where \(\Pi_l\) cycles the selected axes and
\(s_{l,j}>0\) is a scalar.

The geometry and the growth argument use different features of the same
schedule. In period $j$, the control intervals have length comparable to
$j^6$ and the nonround dwell has length $j^7$, while the accumulated
logarithmic radius is comparable to $j^8$. Angular motion is consequently
slow enough for a small concave radial tail to preserve nonnegative Ricci
curvature. The diagonal leg contributes
$\log(D_{l,j})_{ii}=j^7(d_{l,i}+o(1))$, while
$\log s_{l,j}=O(j^6)$. Exact cycling and the asymptotic equality of finitely
many consecutive dwell lengths yield the means in
\eqref{eq:overview-selection}. Bounded Riccati matrices control the intervals
between period endpoints. Since a whole period is negligible compared with
the accumulated logarithmic radius, the same exponent bound holds at every
radius. The strict inequality below $k$ then gives the pointwise growth
condition, and invertibility of the regular fundamental matrices gives the
claimed number of independent smooth harmonic functions.

All spaces and functions are real unless explicitly complexified. Angular
Laplacians are nonnegative. Matrices act on column vectors, and positivity of
a matrix means positivity as a symmetric operator. Constants may depend on
the fixed finite degree range and fixed compact family of controls; they are
uniform in the period index and in the preceding admissible controls.


Sections~\ref{sec:links} and \ref{sec:counting} establish the Berger formulas
and the finite spectral surplus. Section~\ref{sec:control} proves simultaneous
real control, and Section~\ref{sec:geometry} realizes every admitted control
history by a complete metric with nonnegative Ricci curvature.
Section~\ref{sec:transmission} tunes the full-period transmissions exactly;
Section~\ref{sec:conclusion} proves the every-point growth estimate and the
cone conclusions. Appendix~\ref{sec:finite-selection} records the independent
finite arithmetic selection.

\section{Volume-normalized Berger links}
\label{sec:links}

We need angular metrics that retain fixed finite-dimensional harmonic spaces
while varying their spectra. This section derives their curvature and angular
operators; the next section selects the directions that will be cycled.

Fix an odd integer \(m=2s-1\), with \(s\geq4\), and let \(g_0\) be the
unit round metric on \(S^m\subset\R^{2s}\).
An orthogonal complex structure is a real linear map \(J\) satisfying
\(J^2=-\Id\) and \(J^{\mathsf T}J=\Id\).
It determines the unit round Killing field and its dual one-form
\[
 \xi_J(x)=Jx,\qquad \eta_J=g_0(\xi_J,\cdot).
\]
Write
\[
 P_J=\xi_J\otimes\eta_J,\qquad D_Ju=du(\xi_J).
\]
Thus \(P_J\) is the projection onto the Hopf direction.
For \(a,q>0\), define
\[
 \widetilde g_{q,J}=g_0+(q-1)\eta_J^2,\qquad
 G(a,q,J)=a^2q^{-1/m}\widetilde g_{q,J}.
\]
The factor \(q^{-1/m}\) makes the volume independent of \(q\).
All orthogonal complex structures will be allowed; we do not restrict
them to one component of their parameter space.

This volume normalization and the corresponding Laplacian decomposition
are classical; see Tanno~\cite[Sections~2--4]{Tanno}.
His parameter \(t\) corresponds to \(q=t^m\) when \(a=1\).
Classical curvature calculations for these spheres appear in
Bourguignon--Karcher~\cite[Section~6.6]{BourguignonKarcher}.
We record the formulas and their direct proofs in our parameters below.

For \(l\geq0\), let \(V_l\) be the real vector space of restrictions to
\(S^m\) of homogeneous harmonic polynomials of degree \(l\) on
\(\R^{m+1}\).  Equip \(V_l\) with the round \(L^2\) inner product and put
\[
 N_l=\dim_{\R}V_l
     =\binom{m+l}{m}-\binom{m+l-2}{m},
 \qquad B_l=l(l+m-1),
\]
where the second binomial coefficient is zero for \(l<2\).
For completeness, on homogeneous polynomials give distinct monomials
orthogonal inner products with
\(\langle x^\alpha,x^\alpha\rangle=\alpha!\).
The adjoint of the Euclidean polynomial Laplacian from degree \(l\) to
degree \(l-2\) is multiplication by \(|x|^2\), which is injective.
The Laplacian is therefore surjective, giving the dimension formula.
Restriction to the sphere is injective on homogeneous polynomials.
Euclidean polar coordinates give
\[
 \Delta_{g_0}^{+}\big|_{V_l}=B_l\Id,
 \qquad \Delta_h^{+}:=-\operatorname{div}_h\nabla_h.
\]
Rotations preserve \(V_l\).  In particular, \(D_J\) preserves \(V_l\)
and is skew-adjoint for its round inner product.

\begin{lemma}\label{lem:berger-metrics}
The volume form of \(G(a,q,J)\) is \(a^m\,d\operatorname{vol}_{g_0}\).
Its Ricci endomorphism has eigenvalues
\[
 a^{-2}q^{1/m}\bigl((m-1)-2(q-1)\bigr)
 \quad\hbox{horizontally},\qquad
 a^{-2}q^{1/m}(m-1)q
 \quad\hbox{vertically}.
\]
Consequently,
\begin{equation}\label{eq:link-margin}
 \Ric_{G(a,q,J)}>(m-1)G(a,q,J)
 \quad\Longleftrightarrow\quad
 a^2<q^{1/m}
       \min\left\{1-\frac{2(q-1)}{m-1},\,q\right\}.
\end{equation}
Every \(V_l\) is preserved by the positive Laplacian of \(G(a,q,J)\),
whose restriction is
\begin{equation}\label{eq:angular-operator}
 A_l(a,q,J)
 =a^{-2}q^{1/m}
       \left[B_l\Id+(q^{-1}-1)(-D_J^2|_{V_l})\right].
\end{equation}
\end{lemma}

\begin{proof}
We suppress \(J\) from the notation.
Relative to \(g_0\), the metric \(\widetilde g_q\) has eigenvalues \(1\)
on the \((m-1)\)-dimensional horizontal space and \(q\) on the Hopf
direction.  Its volume form is \(\sqrt q\,d\operatorname{vol}_{g_0}\).
Multiplication of the metric by \(a^2q^{-1/m}\) multiplies this volume
form by \(a^m q^{-1/2}\), proving the volume assertion.

Here is a direct curvature calculation.
Let \(\nabla\) be the round connection and set
\(\phi X=\nabla_X\xi\).  Tangent projection of the ambient derivative
gives
\[
 \phi X=JX+\eta(X)x,\qquad
 \phi^*=-\phi,\qquad \phi\xi=0,\qquad
 \phi^2=-\Id+\xi\otimes\eta,
\]
and differentiation gives
\[
 (\nabla_X\eta)(Y)=g_0(\phi X,Y),\qquad
 (\nabla_X\phi)Y=\eta(Y)X-g_0(X,Y)\xi.
\]
Put \(\delta=q-1\), which is constant on the sphere.
If \(C=\widetilde\nabla-\nabla\) is the difference between the
connections of \(\widetilde g_q\) and \(g_0\), the Koszul formula yields
\begin{align*}
 2\widetilde g_q(C(X,Y),Z)
 &=(\nabla_X\widetilde g_q)(Y,Z)
   +(\nabla_Y\widetilde g_q)(X,Z)
   -(\nabla_Z\widetilde g_q)(X,Y)\\
 &=2\delta\left[\eta(X)g_0(\phi Y,Z)
                    +\eta(Y)g_0(\phi X,Z)\right].
\end{align*}
The vectors \(\phi X,\phi Y\) are horizontal, so
\[
 C(X,Y)=\delta\bigl(\eta(X)\phi Y+\eta(Y)\phi X\bigr).
\]
In particular, \(\tr[X\mapsto C(X,Z)]=0\).
Tracing the connection change formula for curvature therefore gives
\[
 (\Ric_{\widetilde g_q}-\Ric_{g_0})(Y,Z)
 =\tr\bigl[X\mapsto(\nabla_XC)(Y,Z)\bigr]
  -\tr\bigl[X\mapsto C(Y,C(X,Z))\bigr].
\]
These are endomorphism traces and can be computed in a round
orthonormal frame.
Using the identities for \(\phi\), the four terms obtained by
differentiating \(C\) contribute, in order,
\[
 -\delta(g_0-\eta^2),\quad
 \delta(m-1)\eta^2,\quad
 -\delta(g_0-\eta^2),\quad
 \delta(m-1)\eta^2.
\]
Thus the first trace is
\(\delta(-2g_0+2m\eta^2)(Y,Z)\).
For the second trace, one has
\[
 C(Y,C(X,Z))
 =\delta^2\eta(Y)
   \bigl[-\eta(X)Z-\eta(Z)X+2\eta(X)\eta(Z)\xi\bigr],
\]
whose trace in \(X\) is
\(-(m-1)\delta^2\eta(Y)\eta(Z)\).
Since \(\Ric_{g_0}=(m-1)g_0\), this proves the tensor identity
\[
 \Ric_{\widetilde g_q}
 =\bigl((m-1)-2(q-1)\bigr)(g_0-\eta^2)
   +(m-1)q^2\eta^2.
\]
The squared length of \(\xi\) in \(\widetilde g_q\) is \(q\).
Hence the vertical eigenvalue of its Ricci endomorphism is
\((m-1)q\), while the horizontal eigenvalue is
\((m-1)-2(q-1)\).
Constant metric scaling leaves the covariant Ricci tensor unchanged
and divides its endomorphism eigenvalues by the scale.
This proves the stated eigenvalues and \eqref{eq:link-margin}.

Finally, the inverse metric is
\[
 G(a,q,J)^{-1}
 =a^{-2}q^{1/m}
     \left[g_0^{-1}+(q^{-1}-1)\xi\otimes\xi\right].
\]
Because its volume form is a spatially constant multiple of round
volume, its divergence on vector fields is the round divergence.
The field \(\xi\) has zero round divergence, so
\(\operatorname{div}_{g_0}((D_Ju)\xi)=D_J^2u\).
It follows that
\[
 \Delta_{G(a,q,J)}^{+}u
 =a^{-2}q^{1/m}
       \left[\Delta_{g_0}^{+}u+(q^{-1}-1)(-D_J^2u)\right].
\]
Restricting this identity to \(V_l\) gives
\eqref{eq:angular-operator}.
\end{proof}

\section{A strict spectral surplus}\label{sec:counting}

Fix the standard orthogonal complex structure $J_0$ on $\mathbb R^{2s}$.
For $a,q>0$, let
$\lambda_{l,1}(a,q)\le\cdots\le\lambda_{l,N_l}(a,q)$
be the eigenvalues, counted with their real multiplicities, of the positive
angular operator in \eqref{eq:angular-operator} for $G(a,q,J_0)$.
Define the corresponding nonnegative exponents by
\begin{equation}\label{eq:dwell-exponents}
 d_{l,i}(a,q)\bigl(d_{l,i}(a,q)+m-1\bigr)=\lambda_{l,i}(a,q),
 \qquad d_{l,i}(a,q)\ge0.
\end{equation}
We will find more than the Euclidean number of directions whose exponents
have mean strictly below a common integer. The subsequent construction
will realize these means by cycling the selected directions.

The Hopf-charge decomposition underlying this calculation is described
by Tanno~\cite[Lemmas~3.1 and~4.1]{Tanno}; we prove the exact
multiplicities and the distribution needed for the count.

\begin{lemma}[Hopf-weight distribution]\label{lem:hopf-distribution}
For $m=2s-1$ and $0\le b\le l$, the eigenvalue $(2b-l)^2$ of
$-D_{J_0}^2$ on $V_l$ has the following contribution to its multiplicity:
\begin{align}
 h_{l,b}
 &=\binom{s+b-1}{s-1}\binom{s+l-b-1}{s-1}
   -\binom{s+b-2}{s-1}\binom{s+l-b-2}{s-1} \label{eq:hopf-multiplicity}\\
 &=\frac{l+s-1}{s-1}
   \binom{b+s-2}{s-2}\binom{l-b+s-2}{s-2}.\notag
\end{align}
Terms with a missing bidegree in the first line are zero.
With probabilities $h_{l,b}/N_l$, the quantities $(2b-l)^2/l^2$
converge in distribution, as $l\to\infty$, to
\[
 Y_m\sim\operatorname{Beta}\left(\frac12,\frac{m-1}{2}\right).
\]
\end{lemma}

\begin{proof}
Complexify $V_l$ and decompose homogeneous polynomials into bidegrees
$(b,l-b)$ in $z_1,\ldots,z_s$ and their conjugates.
The operator $D_{J_0}$ acts on this bidegree by $\mathrm i(2b-l)$.
The contraction $\sum_j\partial_{z_j}\partial_{\bar z_j}$ maps onto
the polynomials of bidegree $(b-1,l-b-1)$: in the positive monomial
inner product in which $z^\alpha\bar z^\beta$ has squared norm
$\alpha!\beta!$, its adjoint is multiplication by
$\sum_j z_j\bar z_j$, which is injective.
Taking the difference of dimensions gives the first line of
\eqref{eq:hopf-multiplicity}; cancellation gives the second.
The contributions indexed by \(b\) and \(l-b\) belong to the same
eigenvalue of \(-D_{J_0}^2\) and are added; the middle contribution is counted
once. Since this operator is real symmetric, the complex dimension of each
complexified eigenspace equals the real dimension of the corresponding
original real eigenspace. There is no additional
factor of two.

Summing the second line by the binomial convolution identity gives
\begin{equation}\label{eq:harmonic-multiplicity-asymptotic}
 N_l=\frac{l+s-1}{s-1}\binom{l+2s-3}{2s-3}
 =\binom{m+l}{m}-\binom{m+l-2}{m}
 \sim\frac{2}{(m-1)!}\,l^{m-1}.
\end{equation}
More precisely, $h_{l,b}/N_l$ is a beta-binomial probability.
If $X$ has density proportional to $[x(1-x)]^{s-2}$ on $(0,1)$ and,
conditionally on $X$, $K_l$ is binomial with parameters $l,X$, then
integration of the binomial probability against this density gives
$\mathbb P(K_l=b)=h_{l,b}/N_l$.
Since
\[
 \mathbb E\left[\left(\frac{K_l}{l}-X\right)^2\right]
 =\frac{\mathbb E[X(1-X)]}{l}\le\frac{1}{4l},
\]
we have $K_l/l\to X$ in probability in this representation.
The change of variable $y=(2x-1)^2$ gives density proportional to
$y^{-1/2}(1-y)^{s-2}$, proving the assertion.
\end{proof}

For $0<v<1$, define the upper-tail quantile $y_m(v)\in(0,1)$ by
\[
 \mathbb P\bigl(Y_m>y_m(v)\bigr)=v,
 \qquad
 \mathcal J_m:=\int_0^1 y_m(v)\log(1/v)\,dv.
\]
The threshold $2/(m-1)$ in the next lemma comes from the horizontal
factor in the link curvature condition~\eqref{eq:link-margin}.

\begin{lemma}[A weighted quantile inequality]\label{lem:weighted-quantile}
For all sufficiently large odd $m$,
\begin{equation}\label{eq:weighted-quantile-gap}
 (m-1)\mathcal J_m>2.
\end{equation}
\end{lemma}

\begin{proof}
Put $r=m-1$. The beta density, equivalently polar coordinates for
independent standard normal variables, gives the representation
\[
 rY_m\overset{\mathrm d}=\frac{Z^2}{Z^2/r+W_r/r},
\]
where $Z$ is standard normal and $W_r$ is an independent sum of $r$
squared standard normal variables. Since
$\mathbb E[W_r/r]=1$ and $\operatorname{Var}(W_r/r)=2/r$,
we obtain $rY_m\Rightarrow Z^2$.
The limiting distribution is continuous and strictly increasing on
$(0,\infty)$, so its upper quantile $y_\infty(v)$ satisfies
$r y_m(v)\to y_\infty(v)$ for every $0<v<1$.
Fatou's lemma therefore gives
\begin{equation}\label{eq:normal-quantile-limit}
 \liminf_{m\to\infty}(m-1)\mathcal J_m
 \ge \int_0^1 y_\infty(v)\log(1/v)\,dv
 =\mathbb E\left[R^2\log\frac{1}{Q(R)}\right],
\end{equation}
where $R=|Z|$ and $Q(x)=\mathbb P(R>x)$.
The last equality follows since $Q(R)$ is uniform on $(0,1)$.
This expectation is finite: integration of the half-normal density on
$[x,x+1]$ gives
$Q(x)\ge\sqrt{2/\pi}\,e^{-(x+1)^2/2}$.

For $x>0$, integration of the same density gives the classical Mills
upper bound (see \cite[Section~4, p.~113]{Lee1992})
\[
 Q(x)=\sqrt{\frac2\pi}\int_x^\infty e^{-t^2/2}\,dt
 \le\sqrt{\frac2\pi}\,\frac{e^{-x^2/2}}{x}.
\]
Consequently, writing $A=\mathbb E[R^2\log R]$, we have
\begin{equation}\label{eq:normal-quantile-lower-bound}
 \mathbb E\left[R^2\log\frac1{Q(R)}\right]
 \ge\frac32+A+\frac12\log\frac\pi2.
\end{equation}
Integration by parts against the half-normal density yields
$\mathbb E[R^4\log R]=3A+1$.
Under the probability measure $d\nu=R^4\,d\mathbb P/3$, Jensen's
inequality for the logarithm applied to $1/R$ now gives
\[
 A+\frac13=\mathbb E_\nu[\log R]
 \ge-\log\mathbb E_\nu[1/R]
 =\log\frac3{\mathbb E[R^3]}
 =\log\frac3{2\sqrt{2/\pi}}.
\]
Combining this with \eqref{eq:normal-quantile-lower-bound} proves
\begin{equation}\label{eq:strict-normal-constant}
 \mathbb E\left[R^2\log\frac1{Q(R)}\right]
 \ge\frac76+\log\frac{3\pi}{4}>2.
\end{equation}
Here the final strict inequality can be checked without decimal
approximations. Indeed, $\arctan x>x-x^3/3$ for $x>0$ gives
\[
 \frac\pi4=\arctan\frac12+\arctan\frac13
 >\frac{505}{648}>\frac79,
\]
so $3\pi/4>7/3$. Strict convexity of $t\mapsto1/t$ gives the
midpoint bounds
\[
 \log\frac73
 =\int_1^{5/3}\frac{dt}{t}+\int_{5/3}^{7/3}\frac{dt}{t}
 >\frac{2/3}{4/3}+\frac{2/3}{2}=\frac56.
\]
Equations~\eqref{eq:normal-quantile-limit} and
\eqref{eq:strict-normal-constant} imply \eqref{eq:weighted-quantile-gap}.
\end{proof}

Before selecting the link, fix a smooth nondecreasing cutoff \(\chi\),
zero near \(( -\infty,2]\) and one near \([4,\infty)\), and put
\[
 I_\chi=\int_2^\infty\chi(t)(1+t)^{-5/4}\,dt>0.
\]
The radial construction will require only the extra smallness condition
\begin{equation}\label{eq:cutoff-smallness}
 \frac{-\log a_*}{I_\chi}\,
 \sup_{t>2}\chi(t)(1+t)^{-5/4}\leq\frac14.
\end{equation}
It depends on the fixed cutoff and link scale, before any finite degree
range or control family has been chosen. The freedom in the next lemma
allows this condition to be imposed at the outset.

\begin{lemma}[Spectral surplus]\label{lem:spectral-surplus}
There is an odd $m=2s-1\ge7$ for which the following choices are possible.
The parameters $a_*\in(0,1)$ and $q_*>1$ can be taken arbitrarily close
to $1$, with \eqref{eq:link-margin} holding at $a=a_*$ for every
$q\in[1,q_*]$.
For these parameters there are integers $k\ge2$, $L\ge2$, and
$0\le M_l\le N_l$ for $2\le l\le L$ such that, on abbreviating
$d_{l,i}=d_{l,i}(a_*,q_*)$,
\begin{equation}\label{eq:selection}
 \overline d_l:=\frac1{M_l}\sum_{i=1}^{M_l}d_{l,i}<k
 \quad\text{whenever }M_l>0,
 \qquad
 \sum_{l=2}^L M_l>\sum_{l=0}^k N_l.
\end{equation}
\end{lemma}

\begin{proof}
\noindent\textbf{Choosing the link.}
Choose an odd $m\ge7$ satisfying Lemma~\ref{lem:weighted-quantile}, and
fix
\[
 \frac2{m-1}<c<\mathcal J_m.
\]
For a small positive parameter $\varepsilon$, set
\begin{equation}\label{eq:dwell-parameters}
 q_*=1+\varepsilon,
 \qquad \alpha_*^2=1-c\varepsilon,
 \qquad a_*^2=q_*^{1/m}\alpha_*^2.
\end{equation}
For $q\ge1$, the right-hand side of \eqref{eq:link-margin} is
\[
 F(q)=q^{1/m}\left(1-\frac{2(q-1)}{m-1}\right),
 \qquad
 F'(q)=\frac{m+1}{m(m-1)}q^{1/m-1}(1-2q)<0.
\]
At $q=q_*$, the required strict inequality follows from
$c>2/(m-1)$; it consequently holds throughout $[1,q_*]$.
In particular $a_*^2<F(q_*)<1$.

\medskip\noindent\textbf{Limiting means.}
By \eqref{eq:angular-operator}, the eigenvalues at the dwell parameters
are the numbers
\[
 \alpha_*^{-2}\left[B_l-
       \left(1-\frac1{q_*}\right)(2b-l)^2\right]
\]
with the multiplicities in Lemma~\ref{lem:hopf-distribution}.
Thus their increasing order corresponds to the decreasing order of the
Hopf weights $(2b-l)^2$.
If $d(d+m-1)=\lambda$ and $d\ge0$, then
$0\le\sqrt\lambda-d\le(m-1)/2$.
Since $B_l/l^2\to1$, Lemma~\ref{lem:hopf-distribution} and convergence
of quantiles to a continuous, strictly increasing distribution imply
\begin{equation}\label{eq:limiting-exponent-quantile}
 \frac{d_{l,\lceil vN_l\rceil}}{l}\longrightarrow
 f_\varepsilon(v):=(1-c\varepsilon)^{-1/2}
 \sqrt{1-\left(1-\frac1{q_*}\right)y_m(v)}
 \qquad(0<v<1).
\end{equation}
These rescaled exponents are uniformly bounded, so averaging the first
$\lceil uN_l\rceil$ terms gives, for every $0<u<1$,
\begin{equation}\label{eq:limiting-mean-exponents}
 \frac1{l\lceil uN_l\rceil}
       \sum_{i=1}^{\lceil uN_l\rceil}d_{l,i}
 \longrightarrow
 g_\varepsilon(u):=\frac1u\int_0^u f_\varepsilon(v)\,dv.
\end{equation}
One may see this directly by integrating the empirical quantile in
\eqref{eq:limiting-exponent-quantile}; the possible last fractional term
has size $O(N_l^{-1})$ after division by $l$.

Define
\[
 I(\varepsilon):=\int_0^1 g_\varepsilon(u)^{-m}\,du.
\]
All functions here are bounded away from zero for sufficiently small
$\varepsilon$. Uniformly for $0<v<1$,
\[
 f_\varepsilon(v)
 =1+\frac\varepsilon2\bigl(c-y_m(v)\bigr)+O(\varepsilon^2).
\]
The same uniform remainder remains valid after averaging over $(0,u)$.
Expansion of $g_\varepsilon(u)^{-m}$ and integration therefore give
\begin{align}
 I(\varepsilon)
 &=1+\frac{m\varepsilon}{2}
   \left[\int_0^1\frac1u\int_0^u y_m(v)\,dv\,du-c\right]
   +O(\varepsilon^2)\notag\\
 &=1+\frac{m\varepsilon}{2}(\mathcal J_m-c)
   +O(\varepsilon^2)>1
 \qquad\text{for all sufficiently small }\varepsilon>0.
 \label{eq:counting-first-variation}
\end{align}
The second equality uses the nonnegative integral identity
$\int_v^1du/u=\log(1/v)$.
Fix such an $\varepsilon$, as small as desired, from now on.

\medskip\noindent\textbf{Counting at integer thresholds.}
For each positive integer $k$ and $l\ge2$, let $M_l(k)$ be the largest
initial group of the ordered exponents whose mean is strictly less than
$k$, taking $M_l(k)=0$ if no group qualifies.
Only finitely many degrees contribute: the displayed eigenvalue formula
bounds all $\lambda_{l,i}$ below by a positive constant times $l^2$,
so $\min_i d_{l,i}$ grows at least linearly in $l$.
The initial-group means are nondecreasing, and hence
\begin{equation}\label{eq:counting-layer-cake}
 M_l(k)=N_l\int_0^1
 \mathbf 1\left\{
   \frac1{\lceil uN_l\rceil}
   \sum_{i=1}^{\lceil uN_l\rceil}d_{l,i}<k
 \right\}\,du.
\end{equation}
For fixed $u\in(0,1)$ and $\delta>0$,
\eqref{eq:limiting-mean-exponents} bounds the mean in this indicator by
$l(g_\varepsilon(u)+\delta/2)$ for every sufficiently large $l$.
It is therefore strictly below $k$ whenever
$l\le k/(g_\varepsilon(u)+\delta)$, including at the upper endpoint. Using
\eqref{eq:harmonic-multiplicity-asymptotic}, the sum of the corresponding
$N_l$, divided by $\sum_{l=0}^kN_l$, has limit
$(g_\varepsilon(u)+\delta)^{-m}$ as the integer $k\to\infty$.
Letting $\delta\downarrow0$ and applying Fatou's lemma to
\eqref{eq:counting-layer-cake} yields
\begin{equation}\label{eq:counting-liminf}
 \liminf_{\substack{k\to\infty\\ k\in\mathbb N}}
 \frac{\sum_{l=2}^\infty M_l(k)}{\sum_{l=0}^kN_l}
 \ge\int_0^1g_\varepsilon(u)^{-m}\,du
 =I(\varepsilon)>1.
\end{equation}
Discarding the finitely many lower degrees in this argument changes no
limit. Consequently every sufficiently large integer $k$ has a strict surplus.
Taking $L$ to be the largest degree with $M_l(k)>0$ and setting
$M_l=M_l(k)$ proves \eqref{eq:selection}.
\end{proof}

A selected prefix may be empty, a singleton, or all of \(V_l\), and may
cut a repeated real eigenspace. No gap at its last eigenvalue is required:
choose the indicated number of real orthonormal eigenvectors. The selected
span need not be invariant under any individual \(D_J\); the full fixed
space \(V_l\) remains invariant under every angular operator.

Finally, the sum appearing in \eqref{eq:selection} is exactly the
Euclidean comparison dimension in ambient dimension $n=m+1$:
\[
 \sum_{l=0}^kN_l
 =\binom{m+k}{k}+\binom{m+k-1}{k-1}
 =h_k(\mathbb R^{m+1}).
\]
To justify the Euclidean identification directly, let $u$ be a Euclidean
harmonic function of growth at most the integer $k$. A smooth radial kernel
$\psi_R(x)=R^{-(m+1)}\psi(x/R)$, supported in the ball of radius $R$ and
with integral $1$, reproduces $u$ by the mean-value property. For each
multi-index $\alpha$ of order $k+1$, moving derivatives onto the kernel
gives, at any fixed $x$ and for $R\ge1$,
\[
 |\partial^\alpha u(x)|
 \le \sup_{|y-x|\le R}|u(y)|\,
       \|\partial^\alpha\psi_R\|_{L^1}
 \le C_x R^kR^{-(k+1)}\longrightarrow0.
\]
Thus $u$ is a polynomial of degree at most $k$, and its homogeneous
components are harmonic. The converse follows from homogeneity.

For the rest of the construction, the spectral data mean an odd
\(m=2s-1\) with \(s\geq4\), a fixed cutoff \(\chi\), and
\(0<a_*<1<q_*\) satisfying \eqref{eq:link-margin} throughout
\([1,q_*]\) and \eqref{eq:cutoff-smallness}, together with finite
\(k,L\geq2\) and \(0\leq M_l\leq N_l\) satisfying
\eqref{eq:selection}. These are precisely the numerical hypotheses used
in Sections~\ref{sec:control}--\ref{sec:conclusion}; no later step uses
the asymptotic route by which the data were obtained.

The choices in Lemma~\ref{lem:spectral-surplus} are made in the order
$m,c,\varepsilon,k,L$. In particular any further requirement that $a_*$
be close to $1$ can be imposed when choosing $\varepsilon$, before the
integer cutoff is fixed.

\section{Control of the angular modes}
\label{sec:control}

The selections in Lemma~\ref{lem:spectral-surplus} will be mixed by
cycling their basis directions.  We first prove that the angular
operators furnish the required finite-dimensional control.  Throughout
this section, \(m=2s-1\), \(s\geq4\), and \(2\leq l\leq L\).  We equip
each real space \(V_l\) of round spherical harmonics with its round
\(L^2\) inner product and write \(N_l=\dim_{\mathbb R}V_l\).
For an endomorphism \(A\) of \(V_l\), put
\[
 A_{\mathrm{tf}}=A-\frac{\operatorname{tr}A}{N_l}I.
\]
For every orthogonal complex structure \(J\) on \(\mathbb R^{2s}\), the
operator \(D_Ju(x)=du_x(Jx)\) preserves \(V_l\).  It is skew-adjoint,
because its flow consists of round isometries.  In particular,
\((D_J^2)_{\mathrm{tf}}\) is a real symmetric trace-free operator.

\begin{proposition}[Simultaneous control]
\label{prop:control}
Fix positive numbers \(\kappa_l\), \(2\leq l\leq L\).  As \(J\) ranges
over all orthogonal complex structures on \(\mathbb R^{2s}\), the tuples
\begin{equation}
 X_J=\bigl(\kappa_l(D_J^2|_{V_l})_{\mathrm{tf}}\bigr)_{l=2}^L
 \label{eq:control-generators}
\end{equation}
generate, under real linear combinations and brackets, the full Lie
algebra
\[
 \bigoplus_{l=2}^L\mathfrak{sl}(V_l).
\]
\end{proposition}

\begin{proof}
We first prove the assertion in a single degree \(l\).  The nonzero
factor \(\kappa_l\) has no effect on this assertion.  Let
\(\mathfrak a\subset\mathfrak{sl}(V_l)\) be the real Lie algebra
generated by \((D_J^2)_{\mathrm{tf}}\), and set
\[
 W=V_l\otimes_{\mathbb R}\mathbb C,\qquad
 \mathfrak a_{\mathbb C}=\mathfrak a\otimes_{\mathbb R}\mathbb C.
\]
Transpose on \(W\) will always refer to the complex bilinear extension
of the real round inner product.  The single-degree argument first
produces a rank-one symmetric endomorphism of \(W\).  Round conjugation
and brackets then produce all trace-free endomorphisms of \(W\).
Finally, we separate the different degrees to obtain simultaneous
control on their direct sum.

\medskip
\noindent\textbf{Diagonal operators and individual weight blocks.}
Choose orthonormal coordinates
\((x_1,y_1,\ldots,x_s,y_s)\) and write \(z_i=x_i+\mathrm i y_i\).
For \(\theta\in\mathbb R^s\), let \(R_\theta\) rotate coordinate
plane \(i\) by \(\theta_i\), and let \(T_\theta u(x)=u(R_\theta x)\).
The \emph{weight space} of \(\nu\in\mathbb Z^s\) is the character
eigenspace
\[
 W_\nu=\{u\in W:T_\theta u=e^{\mathrm i\sum_i\nu_i\theta_i}u
                    \text{ for every }\theta\}.
\]
Thus \(z_i\) has weight \(e_i\), \(\bar z_i\) has weight \(-e_i\),
and \(D_{J_0}\) acts on \(W_\nu\) by \(\mathrm i\sum_i\nu_i\).
In particular, a monomial of holomorphic degree \(b\) and
antiholomorphic degree \(l-b\) has Hopf charge \(2b-l\), consistently
with Section~\ref{sec:counting}.  Every occurring weight satisfies
\[
 \sum_{i=1}^s|\nu_i|\leq l,\qquad
 l-\sum_{i=1}^s|\nu_i|\in2\mathbb Z.
\]
Every weight on the extreme layer
\[
 \Lambda_l=\left\{\nu\in\mathbb Z^s:
                    \sum_{i=1}^s|\nu_i|=l\right\}
\]
occurs and has a one-dimensional weight space: it is spanned by the
monomial using \(z_i^{\nu_i}\) when \(\nu_i\geq0\) and
\(\bar z_i^{-\nu_i}\) when \(\nu_i<0\).  This monomial is harmonic,
since in each coordinate plane it involves at most one of \(z_i,\bar
z_i\).  It is the only degree-\(l\) monomial of that weight.

More generally, every weight satisfying the displayed bound and parity
condition occurs.  To check its multiplicity, set
\(a=(l-\sum_i|\nu_i|)/2\), a nonnegative integer.
Every monomial of weight \(\nu\) is the corresponding signed monomial
of degree \(\sum_i|\nu_i|\), multiplied by
\(\prod_i(z_i\bar z_i)^{a_i}\) with \(a_i\ge0\) and
\(\sum_i a_i=a\).  The full polynomial weight space therefore has
dimension \(\binom{a+s-1}{s-1}\).
The polynomial Laplacian preserves the weight and maps onto the
corresponding weight space two degrees lower: its adjoint in the
positive monomial inner product is a nonzero scalar multiple of
multiplication by \(\sum_i z_i\bar z_i\), which preserves weight and
is injective.  Thus the harmonic weight multiplicity is
\[
 \binom{a+s-1}{s-1}-\binom{a+s-2}{s-1}
 =\binom{a+s-2}{s-2},
\]
where the second term is zero when \(a=0\).
This standard weight formula also appears in
Lauret--Miatello--Rossetti~\cite[Lemma~3.6]{LauretMiatelloRossetti}.
In particular, interior weight spaces may have dimension greater than
one.  The interpolation below isolates an entire source-to-target
weight block; we use a rank-one conclusion only when both endpoints
are on the extreme layer.

Conjugating by a round rotation sends the family of generators to
itself.  Thus \(\mathfrak a_{\mathbb C}\) is invariant under the round
torus and contains each torus Fourier component of any of its
elements.  Indeed, Fourier projection is an integral of scalar
multiples of torus conjugates, and a finite-dimensional vector
subspace is closed.

For \(\sigma=(\sigma_1,\ldots,\sigma_s)\in\{\pm1\}^s\), let \(J_\sigma\)
rotate coordinate plane \(i\) with sign \(\sigma_i\).  On the
\(\nu\)-weight space, \(D_{J_\sigma}^2\) acts by
\(-(\sum_i\sigma_i\nu_i)^2\).  For \(i\ne j\), the identity
\[
 2^{-s}\sum_{\sigma\in\{\pm1\}^s}
 \sigma_i\sigma_j\left[-\left(\sum_a\sigma_a\nu_a\right)^2\right]
 =-2\nu_i\nu_j
\]
shows that \(\mathfrak a_{\mathbb C}\) contains a diagonal operator
\(A_{ij}\) acting on weight \(\nu\) by \(\nu_i\nu_j+c_{ij}\), where
\(c_{ij}\) is independent of \(\nu\).  The scalar term has no effect
on brackets.

An endomorphism Fourier component of \emph{weight shift} \(\delta\)
satisfies \(T_\theta A T_\theta^{-1}
=e^{\mathrm i\sum_i\delta_i\theta_i}A\); equivalently, it maps each
\(W_\nu\) into \(W_{\nu+\delta}\), with absent weight spaces understood
to be zero.  On its block from \(W_\nu\) to \(W_{\nu+\delta}\), the
operator \(\operatorname{ad}A_{ij}\) acts by
\begin{equation}
 \nu_i\delta_j+\nu_j\delta_i+\delta_i\delta_j.
 \label{eq:control-block-eigenvalues}
\end{equation}
If \(\delta\) has at least three nonzero coordinates, these joint
eigenvalues distinguish all source weights.  In fact, equality for
two sources, with difference \(v\), implies
\[
 v_i\delta_j+v_j\delta_i=0\qquad(i\ne j).
\]
On the support of \(\delta\), the ratios \(v_i/\delta_i\) have
pairwise sums zero; three such coordinates force every ratio to be
zero.  Pairing any coordinate outside the support with one in the
support then gives \(v_i=0\) there as well.

There are only finitely many weights.  Polynomials in the commuting
operators \(\operatorname{ad}A_{ij}\) therefore project onto an
individual source block within a fixed shift component.  More
explicitly, for each competing source choose one coordinate of
\eqref{eq:control-block-eigenvalues} that distinguishes it from the
desired source, and multiply the corresponding linear interpolation
factors.  These operations preserve \(\mathfrak a_{\mathbb C}\).
When source and target belong to \(\Lambda_l\), the isolated nonzero
block is a rank-one directed matrix.

\medskip
\noindent\textbf{A rank-one symmetric operator.}
Regard an extreme weight as \(l\) signed tokens, with no opposite
signs at the same coordinate.  We claim that the algebra contains
the directed matrices for every legal move
\[
 \text{two equal-sign tokens at one coordinate}
 \quad\longleftrightarrow\quad
 \text{one token at each of two other coordinates},
\]
where the two latter signs are arbitrary and both endpoints must
remain in \(\Lambda_l\).  Such a shift has three nonzero coordinates.
It remains to exhibit a generator with a nonzero block for the move.

Here is an explicit simultaneous choice of the needed coefficients.
For three distinct coordinate indices \(i,j,k\), let \(X_a,Y_a\)
denote their real coordinate vectors and set
\[
 JX_i=X_j,\qquad JY_i=X_k,\qquad JY_j=Y_k,
\]
with the three reverse images determined by \(J^2=-I\).
Use the standard complex structure on the remaining coordinate
planes.  This defines an orthogonal complex structure.  For signed
variables \(u_a^\epsilon=x_a+\mathrm i\epsilon y_a\),
\(\epsilon\in\{\pm1\}\), its vector field satisfies
\[
 \begin{split}
 D_Ju_i^\epsilon&=-x_j-\mathrm i\epsilon x_k,\\
 D_Ju_j^\beta&=x_i-\mathrm i\beta y_k,\\
 D_Ju_k^\gamma&=y_i+\mathrm i\gamma y_j.
 \end{split}
\]
For every choice of the three signs, the relevant replacement coefficients are
\[
 \begin{array}{c@{\qquad}c@{\qquad}c@{\qquad}c}
 u_i^\epsilon\longmapsto u_j^\beta&
 u_i^\epsilon\longmapsto u_k^\gamma&
 u_j^\beta\longmapsto u_i^\epsilon&
 u_k^\gamma\longmapsto u_i^\epsilon
 \\[3pt]
 -\frac12&-\frac{\mathrm i\epsilon}{2}&
 \frac12&-\frac{\mathrm i\epsilon}{2}
 \end{array}
\]
For completeness, there is no cancellation with other terms of
\(D_J^2\).  In the independent complex linear coordinates
\(z_a,\bar z_a\), the identity \(D_J^2u=-u\) for linear functions gives,
on homogeneous degree-\(l\) polynomials,
\[
 D_J^2p=-lp+
   \sum_{u,v}(D_Ju)(D_Jv)\,\partial_u\partial_vp.
\]
If
\[
 p=(u_i^\epsilon)^a(u_j^\beta)^b(u_k^\gamma)^c M
\]
is an extreme monomial, where \(M\) uses the other coordinate
planes, the coefficient of the split monomial is
\(\mathrm i\epsilon\,a(a-1)/2\), and that of the merge monomial is
\(-\mathrm i\epsilon\,bc/2\).  They are nonzero whenever the indicated
move is legal.  Indeed the exponent deficits of the target force
the differentiated variables to be precisely the two variables
specified by the move.  Since rotations commute with the Euclidean
Laplacian, \(D_J^2p\) is already harmonic; no harmonic projection
changes these coefficients.  Trace subtraction does not affect this
shift.  Fourier projection and the block isolation above now give
the claimed directed matrices.

We next join \(-le_1\) to \(le_1\) by a simple path of these moves.
For even \(l=2h\), use
\[
 -2he_1\longrightarrow
 -(2h-2)e_1+e_2+e_3\longrightarrow\cdots
 \longrightarrow he_2+he_3
 \longrightarrow\cdots\longrightarrow 2he_1,
\]
splitting negative pairs at coordinate \(1\) in the first half and
merging the tokens at \(2,3\) into positive pairs at \(1\) in the
second half.  This includes \(l=2\), for which the path is
\(-2e_1,e_2+e_3,2e_1\).

For odd \(l=2h+1\), \(h\geq1\), both starting signs admit the path
\[
 \pm(2h+1)e_1\longrightarrow\cdots
 \longrightarrow\pm e_1+he_2+he_3
 \longrightarrow(h-1)e_2+he_3+2e_4.
\]
The last step merges the remaining token at \(1\) and one positive
token at \(2\) into two positive tokens at \(4\).  Reverse the
positive-start path and append it to the negative-start path.
The resulting path is simple: the two branches have opposite
nonzero first coordinates until their common endpoint, and no
branch repeats a weight.  For \(l=3\), that endpoint is \(e_3+2e_4\).
All intermediate weights remain extreme.

Choose monomial bases on these one-dimensional weight spaces and
normalize their directed matrices to matrix units.  Nested brackets
along a simple path give its endpoint unit, by
\([E_{c,b},E_{b,a}]=E_{c,a}\) for distinct vertices.  Hence
\(\mathfrak a_{\mathbb C}\) contains the unit from weight \(-le_1\)
to weight \(le_1\).  If \(v=z_1^l\), torus invariance of the bilinear
round pairing shows that \(v\) pairs only with weight \(-le_1\), and
\[
 v^{\mathsf T}\bar v
   =\int_{S^{2s-1}}|z_1|^{2l}\,d\operatorname{vol}_{g_0}>0.
\]
The endpoint unit is consequently a nonzero multiple of
\(vv^{\mathsf T}\).  Notice also that \(v^{\mathsf T}v=0\).

\medskip
\noindent\textbf{From the rank-one operator to all trace-free matrices.}
All round conjugates of \(vv^{\mathsf T}\) belong to the algebra.
We first verify that the round orbit of \(v\) spans \(W\).
Up to a nonzero scalar, its members are the polynomials
\((a\cdot x)^l\) with
\[
 a\in\mathbb C^{2s}\setminus\{0\},\qquad a\cdot a=0.
\]
To see this, write \(a=b+\mathrm i c\) with \(b,c\) real.  The
isotropy condition says that \(b,c\) are orthogonal and have equal
positive length.  After rescaling, this ordered pair is the image
of the first coordinate pair under a real special orthogonal
transformation.

On homogeneous degree-\(l\) polynomials use the Fischer bilinear
form
\[
 (p,q)_{\mathrm F}=p(\partial)q
       =\sum_{|\alpha|=l}\alpha!\,p_\alpha q_\alpha.
\]
Its restriction to real harmonic polynomials is positive definite,
so its complex bilinear restriction to \(W\) is nondegenerate.
If \(p\in W\) annihilates the orbit powers, then
\[
 (p,(a\cdot x)^l)_{\mathrm F}=l!\,p(a)
\]
shows that \(p\) vanishes on the isotropic quadric.  Therefore
\(p\) is divisible by the quadratic polynomial
\(Q(x)=\sum_{a=1}^s(x_a^2+y_a^2)\).  An elementary verification
is to divide by this monic polynomial in the last variable \(y_s\);
the remainder has the form \(A(x')y_s+B(x')\), where \(x'\) denotes
the remaining variables.  Evaluating at the two
distinct roots for generic \(x'\) forces both \(A\) and \(B\) to
vanish identically.

Such a harmonic multiple must be zero.  Indeed multiplication by
\(Q\) is adjoint to the Euclidean Laplacian in the
positive Hermitian Fischer form.  If
\(p=Qq\) and \(\Delta p=0\), then
\[
 \|p\|_{\mathrm F}^2
   =\langle Qq,p\rangle_{\mathrm F}
   =\langle q,\Delta p\rangle_{\mathrm F}=0.
\]
Thus no nonzero element annihilates the orbit, proving its span.

For two orbit vectors \(u,w\),
\[
 [uu^{\mathsf T},ww^{\mathsf T}]
  =(u^{\mathsf T}w)
        (uw^{\mathsf T}-wu^{\mathsf T}).
\]
For fixed \(u\ne0\), the factor \(u^{\mathsf T}w\) is not
identically zero on the orbit, by its spanning property and
nondegeneracy of the round pairing.  It is a real-analytic function
of the round rotation.  The real special orthogonal group is
connected, so its nonzero locus is dense.  Division by this factor
and then passage to limits show that
\(\mathfrak a_{\mathbb C}\) contains
\(uw^{\mathsf T}-wu^{\mathsf T}\) for every pair of orbit vectors.
Their span is the full complex skew-symmetric algebra, since the
orbit spans \(W\).

For every real skew-symmetric \(K\), the operator
\(\exp(\operatorname{ad}K)\) preserves \(\mathfrak a_{\mathbb C}\);
this uses only bracket closure in the complex algebra.
It follows that
\(\mathfrak a_{\mathbb C}\) is invariant under conjugation by the
full real group \(SO(V_l)\).  Here one may use that every special
orthogonal matrix is a product of plane rotations.  This group on
the harmonic-function space sends \(v\), up to a scalar, to every
nonzero isotropic vector of \(W\): their real and imaginary parts
are orthogonal equal-length pairs, and \(\dim V_l\geq3\).
Consequently the algebra contains \(ww^{\mathsf T}\) for every
isotropic \(w\in W\).

These matrices span all trace-free symmetric matrices.  In a real
orthonormal basis, the rank-one matrices associated with
\(e_i+\mathrm i e_j\) and \(e_i-\mathrm i e_j\) have sum
\(2(E_{ii}-E_{jj})\) and difference
\(2\mathrm i(E_{ij}+E_{ji})\).  Together with the skew-symmetric
matrices already obtained, they span \(\mathfrak{sl}(W)\).
We have proved
\(\mathfrak a_{\mathbb C}=\mathfrak{sl}(W)\).
Since \(\mathfrak a\subset\mathfrak{sl}(V_l)\) is a real vector
subspace and complexification preserves dimension, it follows
that \(\mathfrak a=\mathfrak{sl}(V_l)\).

\medskip
\noindent\textbf{Independence of the degrees.}
Let \(\mathfrak h\) be the Lie algebra generated by the tuples
\eqref{eq:control-generators}.  Every one-factor projection is full
by what we have just proved.  Moreover
\[
 N_l=\binom{m+l}{m}-\binom{m+l-2}{m}
     =\binom{m+l-1}{m-1}+\binom{m+l-2}{m-1}
\]
is strictly increasing in \(l\).  The real Lie algebras
\(\mathfrak{sl}(V_l)\) are therefore simple and pairwise
nonisomorphic.  Simplicity can be seen directly: in any nonzero
ideal, a bracket with a matrix unit first produces an element with
an off-diagonal entry if necessary; polynomials in the adjoint
action of a generic trace-free diagonal matrix isolate an
off-diagonal unit.  Brackets with the other matrix units then
produce all off-diagonal units and all diagonal differences.

We give the resulting Lie-algebra Goursat argument explicitly; for the
two-factor formulation, see \cite[Section~3]{FigueroaUngureanu2016}.
Inductively,
suppose the projection onto a preceding product \(\mathfrak b\)
is full, and let \(\mathfrak c\) be the next simple factor.
The projected algebra
\(\mathfrak h'\subset\mathfrak b\oplus\mathfrak c\) is surjective
onto both sides.  The kernel of its projection onto \(\mathfrak b\)
is an ideal of \(\mathfrak c\), and hence is either zero or
\(\mathfrak c\).  In the latter case the whole product is present.
In the former case \(\mathfrak h'\) is the graph of a surjective
homomorphism \(\mathfrak b\to\mathfrak c\).
The image of each simple summand of \(\mathfrak b\) is an ideal in
\(\mathfrak c\); a nonzero image would make that summand isomorphic
to \(\mathfrak c\).  This contradicts their distinct dimensions.
The induction proves the claimed full direct sum.
\end{proof}

\begin{remark}
The quantifier over all orthogonal complex structures in
Proposition~\ref{prop:control} includes both orientation classes.
This is used in the averaging over all sign vectors.  In
particular, when \(s=4\), restricting to one class would identify
some complementary sign characters and would not justify that
averaging step.  The explicit complex structure used for the token
moves can itself be completed in either class by reversing it on
an unused coordinate plane.
\end{remark}

For general background on control systems on Lie groups, see
Jurdjevic--Sussmann~\cite{JurdjevicSussmann}.  We give the finite-word
argument with its nonsingular differential directly.

We now pass from the Lie algebra to exact finite products.  For the
integers \(M_l\) in \eqref{eq:selection}, choose the ordered real
orthonormal dwell eigenbasis \(e_{l,1},\ldots,e_{l,N_l}\).
If \(M_l\geq1\), define a signed cycle by
\[
 \Pi_l e_{l,i}=
 \begin{cases}
 e_{l,i+1},&1\leq i<M_l,\\
 (-1)^{M_l-1}e_{l,1},&i=M_l,\\
 e_{l,i},&i>M_l,
 \end{cases}
\]
and set \(\Pi_l=I\) when \(M_l=0\).
The sign of the \(M_l\)-cycle and its extra column sign cancel,
so \(\det\Pi_l=1\).  Thus
\[
 \Pi=(\Pi_l)_{l=2}^L\in
 \mathcal G:=\prod_{l=2}^L SL(V_l).
\]

\begin{corollary}[A finite word with invertible differential]
\label{prop:word}
Let \(d=\sum_{l=2}^L(N_l^2-1)\).
There exist finitely many orthogonal complex structures
\(J_1,\ldots,J_R\), an open neighborhood \(U\) of zero in
\(\mathbb R^d\), and smooth real functions \(c_1,\ldots,c_R\) on
\(U\) such that
\[
 \mathcal W(h)=
   \exp(c_R(h)X_{J_R})\cdots\exp(c_1(h)X_{J_1})
\]
satisfies \(\mathcal W(0)=\Pi\), and
\[
 D\mathcal W(0):\mathbb R^d\longrightarrow T_\Pi\mathcal G
\]
is an isomorphism.  The same assertion holds for any prescribed
target in \(\mathcal G\).
\end{corollary}

\begin{proof}
Let \(\mathcal H\) be the subgroup of \(\mathcal G\) consisting of
finite products of \(\exp(tX_J)\) with arbitrary real \(t\).  Define
\[
 E=\operatorname{span}_{\mathbb R}
       \{\operatorname{Ad}_hX_J:h\in\mathcal H\}.
\]
Conjugation by every generator exponential preserves \(E\).
Differentiating shows that \(E\) is stable under bracketing with
each \(X_J\).  Since \(E\) contains the generators, it contains
their generated Lie algebra, which is the whole Lie algebra of
\(\mathcal G\) by Proposition~\ref{prop:control}.
Choose a basis
\(\operatorname{Ad}_{h_a}X_{K_a}\), \(1\leq a\leq d\), from this
spanning set.  The map
\[
 F(t_1,\ldots,t_d)=
 \bigl(h_d\exp(t_dX_{K_d})h_d^{-1}\bigr)\cdots
 \bigl(h_1\exp(t_1X_{K_1})h_1^{-1}\bigr)
\]
has value \(I\) at zero and differential
\[
 DF(0)(t)=\sum_{a=1}^d t_a\operatorname{Ad}_{h_a}X_{K_a},
\]
which is an isomorphism.  Every factor \(h_a\), and its inverse,
is a fixed finite word of generator exponentials.  Thus every
value of \(F\) is in \(\mathcal H\).

The inverse function theorem gives an identity neighborhood
contained in \(\mathcal H\), making \(\mathcal H\) an open subgroup.
The group \(SL_N(\mathbb R)\) is connected: polar decomposition
reduces this to connectedness of \(SO(N)\) and of the positive
definite determinant-one matrices, the latter being joined to the
identity by \(S^t\), \(0\leq t\leq1\).
Hence \(\mathcal G\) is connected, and its open subgroup
\(\mathcal H\) is all of \(\mathcal G\).

Choose a fixed finite word representing the desired target and
multiply it by \(F\).  The resulting word has that target as its
value at zero and still has invertible differential.  Expanding
the fixed conjugating words gives the asserted form, with fixed
complex structures and smooth real time parameters.
\end{proof}

\begin{remark}
The real times in Corollary~\ref{prop:word} may have either sign.
They will be realized as integrals of bounded smooth bumps \(z\)
in weak metric pulses of the form \(q=1+z/T\).  Both signs are
admissible once \(T\) is sufficiently large, because \(q>0\) and
converges uniformly to \(1\).  Likewise, the class of \(J\) may be
changed while \(q=1\), when the angular metric is independent of
\(J\).  Thus neither the signed times nor the two orientation
classes impose a restriction on the later construction.
\end{remark}

\section{A complete metric with nonnegative Ricci curvature}
\label{sec:geometry}

Fix spectral data in the sense specified at the end of
Section~\ref{sec:counting}.
The geometric properties we need are \(0<a_*<1<q_*\), the strict
inequality \eqref{eq:link-margin} at \(a=a_*\) for every
\(q\in[1,q_*]\), and the fixed cutoff condition
\eqref{eq:cutoff-smallness}.
Fix also the reference complex structure \(J_0\) used for the dwell
metric.
We describe a metric construction that works for any prescribed finite
family of pulse controls.  Its curvature estimates will be uniform
over all subsequent choices of those controls.
Slowly varying links with fixed volume form also appear in
Colding--Naber~\cite[Lemma~2.1]{ColdingNaber}. We give the curvature
calculation for the present family in full.

The radial cutoff and link parameters were fixed before \(k,L\) and
before any pulse family. Condition~\eqref{eq:cutoff-smallness} will give
\(b=a'/a\geq-1/4\) below. Even very large bounds for a fixed finite
pulse family are handled solely by increasing the starting index \(j_0\).
No spectral parameter is changed in response to the controls.

\paragraph{Pulse data and the period schedule.}
Let \(K\) be a compact ball in a finite-dimensional real parameter
space.
Choose finitely many disjoint open subintervals
\(I_1,\ldots,I_R\) of \((0,1)\), with disjoint closures, smooth
functions \(\beta_\nu\) compactly supported in \(I_\nu\), and
orthogonal complex structures \(J_\nu\).
For each \(\nu\), let \(c_\nu(h)\) be smooth on a neighborhood of \(K\).
Set
\[
 z(\tau,h)=\sum_{\nu=1}^R c_\nu(h)\beta_\nu(\tau).
\]
The amplitudes \(c_\nu\) may have either sign.
The functions, complex structures, and compact set \(K\) are fixed
once and for all.
In applications the bumps can be normalized to have integral \(1\),
so their amplitudes specify the flow parameters in
Corollary~\ref{prop:word}.

For \(j\geq1\), define
\begin{equation}\label{eq:periods}
 T_j=j^6,\qquad L_j=j^7,\qquad t_1=10,\qquad
 t_{j+1}=t_j+L_j+4T_j.
\end{equation}
Choose a starting index \(j_0\), to be made large below.
Put \(q(t)=1\) for \(t\leq t_{j_0}\).
For each \(j\geq j_0\), choose a parameter \(h_j\in K\) and divide
\([t_j,t_{j+1}]\) into the following five successive intervals:
\begin{enumerate}
 \item a pulse of duration \(T_j\), on which
 \[
  q(t)=1+T_j^{-1}z\bigl((t-t_j)/T_j,h_j\bigr);
 \]
 on the support of the \(\nu\)-th bump the complex structure is
 \(J_\nu\);
 \item an upward ramp of duration \(T_j\), taking \(q\) from \(1\)
 to \(q_*\);
 \item a dwell interval of duration \(L_j\), on which \(q=q_*\);
 \item a downward ramp of duration \(T_j\), taking \(q\) from
 \(q_*\) to \(1\);
 \item a round rest of duration \(T_j\), on which \(q=1\).
\end{enumerate}
Use \(J_0\) throughout the last four intervals.
Both ramps use a fixed smooth nondecreasing profile
\(\rho:[0,1]\to[0,1]\), equal to \(0\) near \(0\) and \(1\) near
\(1\): their respective formulas in their own rescaled time are
\(1+(q_*-1)\rho\) and \(q_*-(q_*-1)\rho\).
The last four intervals form the diagonal leg of a period, as
illustrated in Figure~\ref{fig:period}.

\begin{figure}[!ht]
\centering
\setlength{\unitlength}{1mm}
\begin{picture}(152,41)
 \put(0,17){\framebox(25,16){\shortstack{\small Pulse\\$q=1+z/T_j$}}}
 \put(25,17){\framebox(24,16){\shortstack{\small Increase\\$1\to q_*$}}}
 \put(49,17){\framebox(55,16){\shortstack{\small Dwell\\$q=q_*$}}}
 \put(104,17){\framebox(24,16){\shortstack{\small Decrease\\$q_*\to1$}}}
 \put(128,17){\framebox(24,16){\shortstack{\small Round rest\\$q=1$}}}
 \put(12.5,12){\makebox(0,0){$T_j$}}
 \put(37,12){\makebox(0,0){$T_j$}}
 \put(76.5,12){\makebox(0,0){$L_j$}}
 \put(116,12){\makebox(0,0){$T_j$}}
 \put(140,12){\makebox(0,0){$T_j$}}
 \put(25,6){\line(1,0){127}}
 \put(25,6){\line(0,1){2}}
 \put(152,6){\line(0,1){2}}
 \put(88.5,2){\makebox(0,0){\small Fixed diagonal leg: $J=J_0$}}
 \put(0,38){\makebox(0,0)[l]{$t_j$}}
 \put(152,38){\makebox(0,0)[r]{$t_{j+1}$}}
\end{picture}
\caption{A period in logarithmic radius, with $T_j=j^6$ and $L_j=j^7$.
The pulse contains finitely many weak bumps with round gaps between them.
The four subsequent pieces form the fixed diagonal leg. Horizontal lengths
are schematic.}
\label{fig:period}
\end{figure}


The notation \(J(t)\) is only a label for the complex structure used
in the metric at time \(t\).  It need not be continuous.
Labels are changed only in open intervals where \(q=1\), and
\(G(a,1,J)=a^2g_0\) is independent of \(J\) there.
The supported bumps provide such round gaps before, between, and
after their supports.

\begin{proposition}\label{prop:metric}
For the fixed spectral data above, there are a smooth function
\(a:\R\to(0,1]\) and an index \(j_{\rm geom}\) such that, for every
\(j_0\geq j_{\rm geom}\), every sequence
\((h_j)_{j\geq j_0}\) in \(K\), used in the preceding scheme, defines
a smooth complete metric on \(\R^{m+1}\) with nonnegative Ricci
curvature:
\[
 g=dr^2+r^2\gamma(\log r),\qquad
 \gamma(t)=G(a(t),q(t),J(t)).
\]
The function \(a\) is identically \(1\) for \(t\leq2\), decreases to
\(a_*\), and, for some \(\mu>0\), satisfies
\begin{equation}\label{eq:scale-tail}
 b(t):=\frac{a'(t)}{a(t)}
       =-\mu(1+t)^{-5/4}\quad(t\geq4),
 \qquad
 \log\frac{a(t)}{a_*}=4\mu(1+t)^{-1/4}\quad(t\geq4).
\end{equation}
The metric is Euclidean near the origin, and its distance from the
origin is exactly \(r\).
Its volume form is
\[
 d\operatorname{vol}_g
   =r^m a(\log r)^m\,dr\,d\operatorname{vol}_{g_0}.
\]
The function \(a\) and the starting index \(j_0\) can be chosen
independently of the sequence of control parameters.
\end{proposition}

\begin{proof}
\textbf{The radial scale.}
Use the fixed smooth nondecreasing cutoff \(\chi\) from
Section~\ref{sec:counting}.
For \(t>2\), put
\[
 b(t)=-\mu\chi(t)(1+t)^{-5/4},\qquad
 a(t)=\exp\left(\int_2^t b(u)\,du\right),
\]
and set \(b=0\), \(a=1\) for \(t\leq2\).
The integral of \(\chi(t)(1+t)^{-5/4}\) over \((2,\infty)\) is
positive and finite.  Thus there is a unique \(\mu>0\) for which
\(a(t)\to a_*\); moreover \(\mu\to0\) as \(a_*\to1\).
Condition~\eqref{eq:cutoff-smallness} gives \(b\geq-1/4\).
Since \(\chi'\geq0\), for \(t>2\) we have
\[
 b'(t)\leq-\frac{5b(t)}{4(1+t)},\qquad
 b+b'+b^2
 \leq b\left(1-\frac{5}{4(1+t)}+b\right)\leq0.
\]
The last factor is positive for \(t>2\) and \(b\geq-1/4\).
The tail identities \eqref{eq:scale-tail} follow by integration.

\medskip\noindent
\textbf{The slice shape and the mixed Ricci tensor.}
Write \(g_r=r^2\gamma(\log r)\) for the metric on a slice.
A dot will denote differentiation in \(t=\log r\).
Identify the tangent bundles of the slices using the product
coordinates.  Their shape endomorphism is
\[
 S=\frac12 g_r^{-1}\partial_r g_r
   =r^{-1}(\Id+B),\qquad
 B=\frac12\gamma^{-1}\dot\gamma=b\Id+H.
\]
Where a nonround part of a bump or a ramp occurs, \(J\) is fixed.
Direct differentiation of \(G\) gives
\begin{equation}\label{eq:shape-anisotropy}
 H=\frac{\dot q}{2q}\left(P_J-\frac1m\Id\right),\qquad
 \tr H=0,\qquad
 \tr(H^2)=\frac{m-1}{4m}\left(\frac{\dot q}{q}\right)^2.
\end{equation}
On round gaps \(H=0\).  Thus the same identities apply everywhere,
without requiring derivatives of the labels \(J(t)\).

The projection \(P_J\) is also the orthogonal projection for
\(\gamma\) onto the Hopf direction.
Indeed, \(\xi_J\) has constant squared length
\(a^2q^{1-1/m}\) on a fixed slice, and its metric dual is the same
constant times \(\eta_J\).
The flow of \(\xi_J\) preserves both \(g_0\) and \(\eta_J\), so
\(\xi_J\) is a Killing field for \(\gamma\).
A Killing field of constant length has zero divergence and
\(\nabla^\gamma_{\xi_J}\xi_J=0\): the latter identity follows by
pairing the Killing equation with \(\xi_J\).
Equivalently, its unit normalization \(E\) satisfies
\[
 \operatorname{div}_\gamma(E\otimes E^\flat)
    =(\operatorname{div}_\gamma E)E^\flat
      +(\nabla^\gamma_EE)^\flat=0.
\]
Therefore \(\operatorname{div}_\gamma P_J=0\), and
\(\operatorname{div}_\gamma H=0\), because the coefficients in
\eqref{eq:shape-anisotropy} are constant on a slice.
The trace \(\tr S=r^{-1}m(1+b)\) is also spatially constant.
The Codazzi formula consequently makes the mixed Ricci tensor
identically zero.

We record the remaining curvature identities explicitly.
For a metric \(dr^2+g_r\), the connection components involving the
radial variable are
\(\Gamma^r_{ij}=-(g_r)_{ik}S^k{}_j\) and
\(\Gamma^i_{rj}=S^i{}_j\).
They give the normal curvature endomorphism
\(-\partial_rS-S^2\).  Its trace is \(\Ric_g(\partial_r,\partial_r)\).
Tracing the tangential Gauss equation contributes
\(S^2-(\tr S)S\) to the intrinsic slice Ricci endomorphism.
Adding the normal contribution gives
\[
 \Ric_g^\#\big|_{TS^m}
   =\Ric_{g_r}^\#-\partial_rS-(\tr S)S.
\]
Here an index of the tangential tensor is raised using \(g_r\).
The mixed formula is
\(\operatorname{div}_{g_r}S-d(\tr S)\), already shown to vanish.
Substitution of \(S=r^{-1}(\Id+B)\), and
\(\tr B=mb\), yields
\begin{equation}\label{eq:radial-slice-ricci}
 \begin{aligned}
  r^2\Ric_g(\partial_r,\partial_r)
     &=-m(b+\dot b+b^2)-\tr(H^2),\\
  r^2\Ric_g^\#\big|_{TS^m}
     &=\Ric_\gamma^\#-\dot B
       -\bigl(m(1+b)-1\bigr)(\Id+B).
 \end{aligned}
\end{equation}
In particular, the \(H^2\) terms cancel from the tangential identity.

\medskip\noindent
\textbf{Uniform positivity on the periods.}
All constants in the estimates that follow may depend on the fixed
pulse data and \(K\), but not on \(j\) or on any choices of \(h_j\).
On a pulse,
\[
 |q-1|=O(T_j^{-1}),\qquad
 |\dot q|=O(T_j^{-2}),\qquad
 |\ddot q|=O(T_j^{-3}).
\]
On either ramp, \(|\dot q|=O(T_j^{-1})\) and
\(|\ddot q|=O(T_j^{-2})\).
There is no variation of \(q\) during the dwell or rest.
Taking \(j_0\) large makes \(q>0\) throughout all pulses, including
those with negative amplitudes.
The metrics on all slices then lie in a fixed uniformly equivalent
family.  Formula \eqref{eq:shape-anisotropy} gives, throughout period
\(j\),
\[
 \|H\|=O(T_j^{-1}),\qquad
 \|\dot H\|=O(T_j^{-2}),\qquad
 \tr(H^2)=O(T_j^{-2})=O(j^{-12}).
\]
The stronger pulse bound is \(\tr(H^2)=O(T_j^{-4})\).
To obtain the derivative estimate, note that \(P_J\) is independent
of \(t\) where \(q\) varies; on every label-changing gap \(H\)
vanishes identically.

From \eqref{eq:periods},
\[
 t_j\sim\frac{j^8}{8},\qquad
 \sup_{t\in[t_j,t_{j+1}]}\left|\frac{t}{t_j}-1\right|\longrightarrow0.
\]
On the scale tail, an exact calculation gives
\[
 -m(b+\dot b+b^2)
 =m\mu(1+t)^{-5/4}
       \left(1-\frac{5}{4(1+t)}-\mu(1+t)^{-5/4}\right).
\]
It has a positive lower bound of order \(j^{-10}\), uniformly on
period \(j\), for all sufficiently large \(j\).
This dominates the \(O(j^{-12})\) term \(\tr(H^2)\).
The radial component in \eqref{eq:radial-slice-ricci} is therefore
positive once \(j_0\) is large enough.

For the tangential component, the strict inequality
\eqref{eq:link-margin} on the compact interval \([1,q_*]\) implies
that, for some fixed \(\delta>0\),
\[
 \Ric_{G(a_*,q,J_0)}^\#\geq(m-1+2\delta)\Id
       \qquad(1\leq q\leq q_*).
\]
All \(J\) give the same intrinsic Ricci eigenvalues.
Since \(a(t)\to a_*\), this estimate persists with \(2\delta\)
replaced by \(\delta\) along sufficiently late ramps, dwells, and
round rests.
The pulse values satisfy \(q\to1\) uniformly; the strict inequality
at \((a_*,1)\) likewise supplies the same conclusion on late pulses,
after decreasing \(\delta\) if necessary.
This argument allows \(q<1\) on a pulse.
Thus all sufficiently late slices satisfy
\[
 \Ric_\gamma^\#\geq(m-1+\delta)\Id.
\]
On the other hand, \(B=b\Id+H\), \(b,\dot b\to0\), and the displayed
bounds for \(H,\dot H\) show uniformly that
\[
 \dot B+\bigl(m(1+b)-1\bigr)(\Id+B)
       =(m-1)\Id+o(1).
\]
The tangential component of
\eqref{eq:radial-slice-ricci} is consequently positive as well.
All the thresholds used here depend only on the fixed data and
\(K\).  Since the mixed component is zero, these estimates prove
\(\Ric_g\geq0\) on every period for every admitted parameter sequence.

\medskip\noindent
\textbf{The initial region and the global manifold.}
Before \(t_{j_0}\) the metric is round in the angular variable:
\[
 g=dr^2+f(r)^2g_0,\qquad f(r)=r\,a(\log r).
\]
The scale construction gives
\[
 f'(r)=a(1+b)\in(0,1],\qquad
 f''(r)=\frac{a}{r}(b+\dot b+b^2)\leq0.
\]
Its radial and tangential sectional curvatures are respectively
\(-f''/f\) and \((1-(f')^2)/f^2\), as follows from the same normal
curvature and Gauss formulas.  They are nonnegative, proving the
required Ricci condition in this region.

The pulse bumps and ramp profiles are constant or flat at all
junctions, and label changes occur only in round open intervals.
Thus \(\gamma\), and hence \(g\), is smooth for \(r>0\).
Only finitely many periods meet any bounded interval of \(t\), so
the infinite sequence of periods introduces no finite accumulation
of junctions.
For \(t\leq2\), equivalently \(r\leq e^2\), both \(a\) and \(q\)
equal \(1\).  The metric is exactly Euclidean there and extends
smoothly across the origin.
The volume formula follows from Lemma~\ref{lem:berger-metrics}.

A radial path from the origin to a point of radius \(r\) has length
\(r\).  Conversely, the length of any such path is at least the
total variation of its radial coordinate, and hence at least \(r\).
Thus \(d_g(0,x)=r(x)\).
The closed centered balls are the compact sets \(\{r\leq R\}\)
in the underlying \(\R^{m+1}\), so the metric is complete.
The underlying manifold is connected, smooth, and without boundary.
\end{proof}

\section{Exact transmission on the regular harmonic subspaces}
\label{sec:transmission}

The link, selection, and control results are now proved. Fix finite spectral
data as specified at the end of Section~\ref{sec:counting};
Lemma~\ref{lem:spectral-surplus} supplies such data. Thus the real spaces
\(V_l\), the finite range \(2\leq l\leq L\), and the selected
multiplicities \(M_l\) are fixed.
Keep the ordered dwell eigenbases and signed cycles \(\Pi_l\) chosen
in Section~\ref{sec:control}. Every matrix in this section acts in those
fixed round orthonormal coordinates.

Our task is to make the harmonic equation realize these cycles after
an entire period, including its long nonround dwell. First we choose the
pulse family using Corollary~\ref{prop:word}. Proposition~\ref{prop:metric}
then supplies the metric for every admitted sequence once the start is
sufficiently late. The remainder of this section proves uniform estimates
for those provisional metrics before choosing the actual sequence.

\subsection{The ideal pulse and its geometric realization}
We now choose the constants in the control proposition to agree with
the response of the radial equation.  Let \(p_\infty=m-1\), and
for \(2\leq l\leq L\) define
\begin{equation}
 \theta_l>0,\qquad
 \theta_l^2+p_\infty\theta_l=a_*^{-2}B_l,\qquad
 \kappa_l=\frac{a_*^{-2}}{p_\infty+2\theta_l}>0.
 \label{eq:round-root}
\end{equation}
Apply Corollary~\ref{prop:word} to the target \(\Pi\) with these \(\kappa_l\).
Restrict its parameter domain to a closed ball \(K\) centered at
zero and contained in \(U\).  Choose open intervals
\(I_1,\ldots,I_R\subset(0,1)\), in chronological order and with
disjoint closures, and smooth functions \(\beta_\nu\) compactly
supported in \(I_\nu\) with
\(\int_0^1\beta_\nu(\tau)\,d\tau=1\).  Put
\begin{align}
 z(\tau,h)&=\sum_{\nu=1}^{R}c_\nu(h)\beta_\nu(\tau),\nonumber\\
 R_l(\tau,h)&=\kappa_l\sum_{\nu=1}^{R}
   c_\nu(h)\beta_\nu(\tau)
   \left(\frac{B_l}{m}I+D_{J_\nu}^2|_{V_l}\right).
 \label{eq:pulse-profile}
\end{align}
These functions, with all derivatives that occur below, are bounded
on the fixed compact parameter set.  Let \(H_l\) solve
\begin{equation}
 \partial_\tau H_l=R_l(\tau,h)H_l,\qquad H_l(0,h)=I.
 \label{eq:ideal-pulse}
\end{equation}
For a matrix \(C\) on \(V_l\) with positive determinant, define
\[
 \mathcal N_l(C)=(\det C)^{-1/N_l}C.
\]
Each bump in \eqref{eq:pulse-profile} is a scalar function of time
times one constant matrix.  It therefore contributes
\(\exp\bigl(c_\nu(h)\kappa_l(B_lI/m+D_{J_\nu}^2)\bigr)\) to
\(H_l(1,h)\), with later bumps multiplying on the left.  For
positive-determinant matrices,
\(\mathcal N_l(C_2C_1)=\mathcal N_l(C_2)\mathcal N_l(C_1)\);
also \(\mathcal N_l(e^C)=e^{C_{\mathrm{tf}}}\) for every real
matrix \(C\).  Thus removing the determinant removes the scalar
part of each exponent.  Consequently
\begin{equation}
 \begin{gathered}
 \Phi(h):=\bigl(\mathcal N_l(H_l(1,h))\bigr)_{l=2}^{L}
          =\mathcal W(h),\\
 \Phi(0)=\Pi,\qquad D\Phi(0)\text{ is an isomorphism}.
 \end{gathered}
 \label{eq:ideal-word}
\end{equation}
The determinants of the \(H_l\) are positive, since they are
exponentials of integrals of real traces.


Apply Proposition~\ref{prop:metric} to this fixed pulse family. Its proof
supplies a threshold \(j_{\rm geom}\) such that every \(j_0\geq j_{\rm geom}\)
and every parameter sequence in \(K\) give a smooth complete metric with
nonnegative Ricci curvature. Use exactly the radial scale \(a\), cutoff,
and five-part schedule \eqref{eq:periods} fixed there. In particular,
\(d_g(0,x)=r(x)\), \(b=a'/a\geq-1/4\), and the angular volume density is
\(a(t)^m\) times round volume. The four pieces after the pulse form the
whole diagonal leg shown in Figure~\ref{fig:period}.

\subsection{The center-regular value equation}
Before analyzing the harmonic equation, we record one estimate
for matrix differences.  It does not require the matrices in a
product to commute.  All matrix norms in its applications are
Frobenius norms from the fixed round inner products; in fixed
finite dimensions they also control operator norms.

\begin{lemma}[A damped matrix difference]
\label{lem:damping}
Let \(E:[u,v]\to\mathbb R^{N\times N}\) be continuously
differentiable, and let \(F,\mathsf L,\mathsf R\) be continuous
matrix-valued functions on the same interval.  Suppose
\[
 E'=F-\mathsf L E-E\mathsf R
\]
there, where \(\mathsf L,\mathsf R\) are
symmetric and
\(\lambda_{\min}(\mathsf L)+\lambda_{\min}(\mathsf R)\geq\gamma>0\).
Then, for \(u\leq t\leq v\),
\begin{align}
 \|E(t)\|&\leq e^{-\gamma(t-u)}\|E(u)\|
       +\int_u^t e^{-\gamma(t-w)}\|F(w)\|\,dw,
       \label{eq:damping-pointwise}\\
 \int_u^v\|E(t)\|\,dt
 &\leq\gamma^{-1}
       \left(\|E(u)\|+\int_u^v\|F(t)\|\,dt\right).
       \label{eq:damping-integrated}
\end{align}
\end{lemma}

\begin{proof}
The Frobenius inner product gives
\[
 \frac12\frac{d}{dt}\|E\|^2
 =\langle E,F\rangle-\langle E,\mathsf L E\rangle
                         -\langle E,E\mathsf R\rangle
 \leq \|E\|\,\|F\|-\gamma\|E\|^2.
\]
The corresponding upper differential inequality for \(\|E\|\)
holds also at its zeros, by continuity (or by first replacing the
norm by \((\|E\|^2+\varepsilon^2)^{1/2}\) and letting
\(\varepsilon\downarrow0\)).  Multiplication by \(e^{\gamma t}\)
and integration gives \eqref{eq:damping-pointwise}.  Integrating
that estimate in \(t\) and interchanging the nonnegative integrals
gives \eqref{eq:damping-integrated}.
\end{proof}

For a harmonic function whose angular part belongs to \(V_l\), write
\[
 u(r,\omega)=\sum_{i=1}^{N_l}Y_i(\log r)e_{l,i}(\omega).
\]
We identify its coefficient vector \(Y(t)\) with
\(\sum_iY_i(t)e_{l,i}\in V_l\).
The function \(u\) is called \emph{center-regular} if it extends smoothly across
the origin.  The volume formula in Proposition~\ref{prop:metric} is the
reason that a closed equation holds on this fixed coefficient
space.  With \(t=\log r\), it is
\begin{equation}
 Y''+p(t)Y'-A_l(t)Y=0,\qquad
 p(t)=m-1+mb(t),\qquad
 A_l(t)=A_l(a(t),q(t),J(t)).
 \label{eq:harmonic-ode}
\end{equation}
Indeed, the radial density is \(r^m a(\log r)^m\).  Its logarithmic
derivative contributes \(m+mb(t)\) in the radial equation, and the
change \(t=\log r\) subtracts \(1\), giving \(p(t)\).
The density is constant in the angular variable, and
\eqref{eq:angular-operator} preserves \(V_l\), so no additional
angular terms or other harmonic degrees occur.

\begin{lemma}[The center-regular value equation]
\label{lem:regular}
There are a threshold \(j_{\rm reg}\geq j_{\rm geom}\) and constants
\(0<c_P<C_P\) such that the following holds for every
\(j_0\geq j_{\rm reg}\), every sequence of parameters in \(K\),
and every \(2\leq l\leq L\).  There is a unique smooth symmetric
matrix \(P_l(t)\) satisfying
\begin{equation}
 P_l'=A_l-pP_l-P_l^2,\qquad P_l=lI\quad(t\leq2).
 \label{eq:riccati}
\end{equation}
It satisfies
\begin{equation}
 c_PI\leq P_l(t)\leq C_PI
 \qquad(t\in\mathbb R,\ 2\leq l\leq L).
 \label{eq:riccati-barriers}
\end{equation}
The solutions of \(Y'=P_lY\) are exactly the center-regular
solutions of \eqref{eq:harmonic-ode}.  Their value vector at
any one time can be prescribed arbitrarily and determines the
solution uniquely.  At each pulse start one also has
\begin{equation}
 \|P_l(t_j)-\theta_l I\|\leq Cj^{-2}
 \qquad(j\geq j_0),
 \label{eq:round-reset}
\end{equation}
including the first pulse.  The constants and the threshold are
uniform over all preceding choices of the parameters in \(K\).
\end{lemma}

\begin{proof}
The matrices \(A_l(t)\) are symmetric.  For sufficiently large
allowed \(j_0\), the values of \(a\) and \(q\) lie in fixed compact
positive intervals, and only finitely many \(V_l\) are involved.
Positivity of \eqref{eq:angular-operator} therefore gives constants
\[
 0<\lambda_-I\leq A_l(t)\leq\lambda_+I.
\]
The radial condition \(b\geq-1/4\), together with \(b\leq0\), gives
\[
 0<p_-:=3m/4-1\leq p(t)\leq p_+:=m-1.
\]
Choose \(c_P>0\) below all the initial values \(l\) so that
\(\lambda_--p_+c_P-c_P^2>0\), and choose \(C_P\) above all of
them so that \(\lambda_+-p_-C_P-C_P^2<0\).  We may decrease
\(c_P\) and increase \(C_P\) so that every \(\theta_l\) also lies
strictly between them.

The Riccati equation preserves symmetry.  At a first contact with
the lower barrier, a unit vector \(v\) with \(P_lv=c_Pv\) satisfies
\[
 v^{\mathsf T}P_l'v
 \geq\lambda_--p_+c_P-c_P^2>0.
\]
At a contact with the upper barrier the corresponding derivative
is at most \(\lambda_+-p_-C_P-C_P^2<0\).  These strict inward
derivatives prevent an eigenvalue from crossing either barrier.
The bounds also keep \(P_l\) in a compact set of matrices on every
finite time interval, so local existence for the ordinary
differential equation continues for all later times.  Uniqueness follows
from the same local ordinary differential equation theorem.  In the
Euclidean region, \(A_l=B_lI\) and \(p=m-1\), so \(P_l=lI\)
indeed satisfies the equation there.

Let \(\mathcal F_l(t)\) be the fundamental matrix of \(Y'=P_lY\)
with \(\mathcal F_l(2)=I\).  Its determinant is
\[
 \det\mathcal F_l(t)
 =\exp\left(\int_2^t\operatorname{tr}P_l(w)\,dw\right)>0.
\]
Moreover,
\[
 Y''=(P_l'+P_l^2)Y=A_lY-pY',
\]
so its columns solve \eqref{eq:harmonic-ode}.  On \(t\leq2\)
they are \(e^{l(t-2)}\) times constant vectors.  Conversely, the
two scalar radial factors for \eqref{eq:harmonic-ode} in this
region are \(e^{lt}\) and \(e^{-(l+m-1)t}\).  Smoothness at the
origin eliminates the second factor.  Thus every center-regular
solution is one of the solutions \(Y'=P_lY\).  Invertibility of
\(\mathcal F_l(t)\) gives the assertion about arbitrary value data.

It remains to prove the reset estimate.  The schedule gives
\begin{equation}
 t_j=\frac{j^8}{8}+O(j^7),\qquad
 \sup_{t\in[t_j,t_{j+1}]}\left|\frac{t}{t_j}-1\right|\longrightarrow0.
 \label{eq:time-asymptotic}
\end{equation}
Times on period \(j\) and on the final rest immediately preceding
it are therefore comparable to \(j^8\).  From
\eqref{eq:scale-tail},
\begin{equation}
 a(t)-a_*=O(j^{-2}),\qquad p(t)-p_\infty=O(j^{-10})
 \label{eq:period-tail}
\end{equation}
uniformly on these intervals.  On a round interval set
\(E_0=P_l-\theta_l I\).  Subtracting \eqref{eq:round-root} from
\eqref{eq:riccati} gives the exact difference equation
\[
 E_0'=f_l^{\mathrm{rd}}(t)I-(pI+P_l)E_0-E_0\theta_l I,\qquad
 f_l^{\mathrm{rd}}(t)=a(t)^{-2}B_l-p(t)\theta_l-\theta_l^2.
\]
Here \(f_l^{\mathrm{rd}}=O(j^{-2})\) by \eqref{eq:period-tail}.  The symmetric
coefficients \(pI+P_l\) and \(\theta_l I\) have a fixed positive
sum of lower eigenvalue bounds.  Lemma~\ref{lem:damping}, on a final
rest of length \(T_{j-1}\), gives
\[
 \|P_l(t_j)-\theta_lI\|
 \leq C e^{-\gamma T_{j-1}}+Cj^{-2}=O(j^{-2}).
\]
The discrepancy at the start of the rest is bounded by
\eqref{eq:riccati-barriers}, independently of the history.  For the
first pulse, the interval up to \(t_{j_0}\) is entirely round.
Apply the same estimate on
\([t_{j_0}/2,t_{j_0}]\) once \(j_0\) is large.  Its length tends
to infinity and all its times are comparable to \(j_0^8\), so the
same bound follows.  This proves \eqref{eq:round-reset} with the
claimed uniformity.
\end{proof}

\subsection{The response of one pulse}
\label{subsec:pulse}

The reset brings the regular radial derivative close to the
round value \(\theta_l\), but its \(O(j^{-2})\) error cannot simply
be multiplied by the pulse duration \(T_j=j^6\).  The next proof
uses an exact scalar radial solution for the long common growth
and integrates the remaining, damped matrix error.  This is the
point at which the denominator \(p_\infty+2\theta_l\) in
\eqref{eq:round-root} enters.  Contraction properties of positive
matrix Riccati flows are classical; see Lawson--Lim~\cite[Theorem~8.5]{LawsonLim}.
We prove the forced estimates, including their parameter
derivatives, in the form used here.

Fix a period \(j\), hold all parameters in earlier periods fixed,
and let the current parameter \(h\) vary in \(K\).  Write
\[
 T=T_j,\qquad \tau=(t-t_j)/T,\qquad
 \eta_j=j^{-2}+T^{-1}.
\]
Let \(\mathcal B_{l,j}(h)\) be the value transmission for
\(Y'=P_lY\) from the start to the end of the pulse, and write
\(P_{l,e}(h)=P_l(t_j+T,h)\).  Define the scalar solution
\begin{equation}
 v_l'=a(t)^{-2}B_l-p(t)v_l-v_l^2,\qquad
 v_l(t_j)=\theta_l,\qquad
 \varphi_{l,j}=\int_{t_j}^{t_j+T}v_l(t)\,dt.
 \label{eq:scalar-baseline}
\end{equation}
This solution is independent of \(h\) and of the earlier
parameters.

\begin{lemma}[Pulse limit with one parameter derivative]
\label{lem:pulse}
For the fixed compact ball and pulse family,
\begin{equation}
 \begin{aligned}
 e^{-\varphi_{l,j}}\mathcal B_{l,j}(h)
     &=H_l(1,h)+O_{C^1}(\eta_j),\\
 P_{l,e}(h)&=\theta_lI+O_{C^1}(\eta_j).
 \end{aligned}
 \label{eq:pulse-limit}
\end{equation}
Here \(O_{C^1}(\eta_j)\) means that the norm and the norm of its
first derivative in the current parameter \(h\), uniformly on
\(K\), are at most \(C\eta_j\).  The constants and the large-\(j\)
threshold are independent of every admitted preceding history.
\end{lemma}

\begin{proof}
The scalar barriers used in Lemma~\ref{lem:regular} also keep \(v_l\)
between fixed positive constants.  Subtracting
\eqref{eq:round-root} from \eqref{eq:scalar-baseline} gives
\[
 (v_l-\theta_l)'=
 \bigl(a(t)^{-2}B_l-p(t)\theta_l-\theta_l^2\bigr)
 -(p+v_l+\theta_l)(v_l-\theta_l).
\]
Its initial value is zero.  The forcing is \(O(j^{-2})\) by
\eqref{eq:period-tail}, and the damping coefficient is uniformly
positive.  Thus
\begin{equation}
 v_l-\theta_l=O(j^{-2})
 \quad\text{throughout the pulse}.
 \label{eq:baseline-close}
\end{equation}
We retain \(v_l\) exactly in the value equation: even the small
pointwise difference in \eqref{eq:baseline-close} can have a large
integral over \(T_j\).

For the angular operator, expand \eqref{eq:angular-operator} at
\(q=1\).  Its first derivative in \(q\) is
\[
 a^{-2}\left(\frac{B_l}{m}I+D_J^2|_{V_l}\right).
\]
In particular the trace-free part has the positive coefficient
\(a^{-2}(D_J^2)_{\mathrm{tf}}\) that appears in
\eqref{eq:control-generators}.  Since \(q=1+z/T\) on the pulse and
\(a(t)-a_*=O(j^{-2})\), Taylor's formula gives
\begin{equation}
 \begin{aligned}
 A_l(t,h)
 &=a(t)^{-2}B_lI+\frac{p_\infty+2\theta_l}{T}R_l(\tau,h)
       +\mathcal E_{A,l}(t,h),\\
 \|\mathcal E_{A,l}\|_{C^1_h}
 &\leq C\left(\frac{j^{-2}}{T}+\frac1{T^2}\right).
 \end{aligned}
 \label{eq:angular-expansion}
\end{equation}
At most one bump is nonzero at any time.  In the gaps the angular
operator is exactly scalar.  Smooth boundedness of the profiles
and their first \(h\)-derivatives proves the uniform remainder
bound, also across those gaps.

Set
\[
 Q_l(t,h)=v_l(t)I+\frac1T R_l(\tau,h).
\]
It is symmetric, and it has a fixed positive lower eigenvalue
bound when \(j\) is large.  Its Riccati residual is
\[
 \rho_l=A_l-pQ_l-Q_l^2-Q_l'.
\]
The scalar terms in this expression cancel by
\eqref{eq:scalar-baseline}.  Using
\(\partial_t(R_l/T)=\partial_\tau R_l/T^2\), the remaining
identity is
\[
 \rho_l
 =-\frac{p-p_\infty+2(v_l-\theta_l)}{T}R_l
   -\frac{R_l^2+\partial_\tau R_l}{T^2}
   +\mathcal E_{A,l}.
\]
Equations \eqref{eq:baseline-close} and
\eqref{eq:angular-expansion} therefore imply
\begin{equation}
 \sup_{t,h}
 \bigl(\|\rho_l\|+\|D_h\rho_l\|\bigr)
 \leq\frac{C\eta_j}{T}.
 \label{eq:residual}
\end{equation}

Let \(E_l=P_l-Q_l\).  No matrices are commuted in the exact
difference equation
\begin{equation}
 E_l'=\rho_l-(pI+P_l)E_l-E_lQ_l.
 \label{eq:pulse-error-equation}
\end{equation}
The lower eigenvalue bounds for \(P_l,Q_l\), and \(p\) allow
Lemma~\ref{lem:damping} with a fixed damping rate.  The profiles
vanish near \(\tau=0\), so
\[
 E_l(t_j)=P_l(t_j)-\theta_lI=O(j^{-2})
\]
by \eqref{eq:round-reset}.  From
\eqref{eq:damping-pointwise}, \eqref{eq:damping-integrated}, and
\eqref{eq:residual}, we obtain
\begin{equation}
 \sup_{t,h}\|E_l(t,h)\|\leq C\eta_j,\qquad
 \sup_h\int_{t_j}^{t_j+T}\|E_l(t,h)\|\,dt\leq C\eta_j.
 \label{eq:error-integrated}
\end{equation}
The initial discrepancy contributes its size to this integral,
because it is damped; it does not contribute its size times \(T\).

We next control the derivative in a current parameter.  Derivatives
are taken with the earlier history fixed.  Write
\(E_{l,h}=D_hE_l\) and \(Q_{l,h}=D_hQ_l\); the calculation applies
to every unit direction in the finite parameter space.
The coefficients of the Riccati equation depend smoothly on \(h\)
on the pulse, so the ordinary differential equation has smooth
parameter dependence there; the uniform barriers keep its solutions
inside one fixed bounded set.
Differentiating \eqref{eq:pulse-error-equation} and using
\(D_hP_l=E_{l,h}+Q_{l,h}\) gives the exact regrouping
\begin{equation}
 E_{l,h}'=D_h\rho_l
 -(pI+P_l)E_{l,h}-E_{l,h}P_l
 -Q_{l,h}E_l-E_lQ_{l,h}.
 \label{eq:parameter-error-equation}
\end{equation}
Its initial value is zero: the past is fixed and
\(R_l(0,h)=D_hR_l(0,h)=0\).  Also
\(\|Q_{l,h}\|\leq C/T\).  Apply Lemma~\ref{lem:damping} to
\eqref{eq:parameter-error-equation}.  Its forcing has pointwise
norm at most \(C\eta_j/T\), by \eqref{eq:residual} and the first
bound in \eqref{eq:error-integrated}.  Its integrated norm is at
most
\[
 C\eta_j+\frac{C}{T}
   \int_{t_j}^{t_j+T}\|E_l\|\,dt
 \leq C\eta_j.
\]
Consequently
\begin{equation}
 \sup_{t,h}\|E_{l,h}\|\leq\frac{C\eta_j}{T},\qquad
 \sup_h\int_{t_j}^{t_j+T}\|E_{l,h}\|\,dt\leq C\eta_j.
 \label{eq:parameter-error-integrated}
\end{equation}
This argument supplies the derivative estimate directly, without
an a priori estimate for \(D_hP_l\).

To pass from coefficient errors to value transmissions, let
\(\mathcal B_l(t,h)\) be the pulse fundamental matrix up to time
\(t\), with value \(I\) at \(t_j\), and set
\[
 F_l(t,h)=
 e^{-\int_{t_j}^t v_l(w)\,dw}\mathcal B_l(t,h),\qquad
 G_l(t,h)=H_l((t-t_j)/T,h).
\]
Then
\[
 F_l'=(R_l/T+E_l)F_l,\qquad
 G_l'=(R_l/T)G_l,\qquad F_l(t_j)=G_l(t_j)=I.
\]
The integrated norms of \(R_l/T\) and \(D_hR_l/T\) are uniformly
bounded; those of \(E_l\) and \(E_{l,h}\) are \(O(\eta_j)\).
The elementary integral inequality
\(u(t)\leq a+\int_{t_j}^t b(w)u(w)\,dw\), \(b\geq0\), implies
\(u(t)\leq a\exp(\int_{t_j}^t b)\): the right-hand side of the
first inequality has derivative at most \(b\) times itself.
Applying this inequality to the fundamental-matrix equations and
their parameter derivatives bounds
\(F_l,G_l,D_hF_l,D_hG_l\) uniformly.

For completeness, \(Z_l=F_l-G_l\) satisfies
\[
 Z_l'=(R_l/T+E_l)Z_l+E_lG_l,\qquad Z_l(t_j)=0.
\]
The integral bound for \(E_l\) gives
\(\sup_t\|Z_l\|\leq C\eta_j\).  Differentiating this equation
introduces the additional terms
\[
 (D_hR_l/T+E_{l,h})Z_l+E_{l,h}G_l+E_lD_hG_l.
\]
Their integrated norms are \(O(\eta_j)\), using the preceding
uniform bounds and \eqref{eq:parameter-error-integrated}.
The same integral inequality gives
\(\sup_t\|D_hZ_l\|\leq C\eta_j\).  Evaluation at the pulse end
proves the first assertion of \eqref{eq:pulse-limit}.
At that end \(R_l(1,h)=D_hR_l(1,h)=0\).  Combining
\(P_l=Q_l+E_l\), \eqref{eq:baseline-close},
\eqref{eq:error-integrated}, and
\eqref{eq:parameter-error-integrated} proves the second assertion.
Every estimate used the earlier history only through the uniform
reset and barriers of Lemma~\ref{lem:regular}, proving the stated
uniformity.
\end{proof}

\subsection{Removing the full diagonal leg}
\label{subsec:leg}

The pulse estimate describes values and radial derivatives at
the pulse end.  The radial derivative need not yet equal
\(\theta_l\) times the value, so diagonal angular coefficients
on the next leg do not by themselves give a diagonal map on
these data.  We now factor the full leg exactly.  This separates
its possibly large unequal growth from the small error in the
entrance derivative.

Let \(\mathcal T_{l,j}(h)\) denote the value transmission of
center-regular solutions from \(t_j\) to \(t_{j+1}\).  On the
diagonal leg let \(U_{l,j}\) and \(V_{l,j}\) be the value
propagation matrices from its start to its end for initial
Cauchy data \((Y,Y')=(I,0)\) and \((0,I)\), respectively.  These
matrices are diagonal in the fixed dwell basis.  Define
\begin{equation}
 D_{l,j}=U_{l,j}+\theta_lV_{l,j},\qquad
 K_{l,j}=D_{l,j}^{-1}V_{l,j}.
 \label{eq:diagonal-reference}
\end{equation}
Thus \(D_{l,j}\) is the value propagation on the entire leg for
the reference data \((I,\theta_l I)\).

\begin{lemma}[Exact normalization of the full leg and dwell growth]
\label{lem:leg}
The matrices \(U_{l,j}\), \(V_{l,j}\), and \(D_{l,j}\) have
positive diagonal entries, and
\begin{equation}
 0\leq K_{l,j}\leq\theta_l^{-1}I.
 \label{eq:K-bound}
\end{equation}
They depend only on the prescribed diagonal leg, not on its
current or earlier pulse parameters.  The exact transmission
identity is
\begin{equation}
 D_{l,j}^{-1}\mathcal T_{l,j}(h)
 =\left[I+K_{l,j}\bigl(P_{l,e}(h)-\theta_lI\bigr)\right]
       \mathcal B_{l,j}(h).
 \label{eq:full-normalization}
\end{equation}
Moreover, uniformly for the finite set of \(l,i\),
\begin{equation}
 \log (D_{l,j})_{ii}
 =L_jd_{l,i}+O(1+L_jj^{-2}+T_j)
 =L_j\bigl(d_{l,i}+o(1)\bigr).
 \label{eq:diagonal-growth}
\end{equation}
\end{lemma}

\begin{proof}
Every operator on the leg is a function of \(I\) and
\(D_{J_0}^2\), and hence is diagonal in the fixed dwell
eigenbasis.  A scalar coordinate satisfies
\[
 y''+p(t)y'-A_i(t)y=0,\qquad A_i(t)>0.
\]
Multiplying the equation for \(y'\) by the positive integrating
factor \(e^{\int p}\) gives
\[
 \bigl(e^{\int p}y'\bigr)'=e^{\int p}A_i y.
\]
Starting with \((y,y')=(1,0)\), positivity of \(y\) makes \(y'\)
nonnegative until any supposed first zero of \(y\); this prevents
that zero.  Starting with \((0,1)\), the solution is positive
for small positive time and the same argument keeps it positive.
Thus the final entries of \(U_{l,j}\) and \(V_{l,j}\) are
positive.  Formula \eqref{eq:diagonal-reference} gives positivity
of \(D_{l,j}\), and entry by entry
\[
 0<\frac{(V_{l,j})_{ii}}
          {(U_{l,j})_{ii}+\theta_l(V_{l,j})_{ii}}
 \leq\theta_l^{-1}.
\]
This proves \eqref{eq:K-bound}.

For an incoming value vector \(y\) at the start of the period,
the data at the pulse end are
\[
 \bigl(\mathcal B_{l,j}y,\,
       P_{l,e}\mathcal B_{l,j}y\bigr).
\]
Their value at the end of the leg is
\[
 \mathcal T_{l,j}y
 =\bigl(U_{l,j}+V_{l,j}P_{l,e}\bigr)\mathcal B_{l,j}y.
\]
Since \(D_{l,j}=U_{l,j}+\theta_lV_{l,j}\), left multiplication
by \(D_{l,j}^{-1}\) gives \eqref{eq:full-normalization}.
The order \(K_{l,j}(P_{l,e}-\theta_lI)\) is part of this
identity; no commutation with \(P_{l,e}\) is used.  The large
matrix \(D_{l,j}\) has disappeared from the error multiplier,
leaving only the uniform bound \eqref{eq:K-bound}.

To estimate \(D_{l,j}\), take its \(i\)-th scalar solution with
initial data \((1,\theta_l)\).  The positivity just proved
allows its logarithmic derivative \(w=y'/y\) throughout the
leg.  It satisfies
\[
 w'=A_{l,i}(t)-p(t)w-w^2.
\]
The scalar barrier argument in Lemma~\ref{lem:regular}, with
initial value \(\theta_l\), keeps \(w\) between fixed positive
constants.  During the dwell,
\[
 A_{l,i}(t)=\lambda_{l,i}+O(j^{-2}),\qquad
 \lambda_{l,i}=d_{l,i}(d_{l,i}+p_\infty),
\]
by \eqref{eq:period-tail} and the fixed value \(q=q_*\).
For \(e=w-d_{l,i}\), subtraction gives
\[
 e'=F_{l,i}(t)-(p+w+d_{l,i})e,\qquad
 F_{l,i}=A_{l,i}-\lambda_{l,i}
                 -(p-p_\infty)d_{l,i}=O(j^{-2}).
\]
The value of \(e\) at the dwell entrance is bounded, and its
damping coefficient has a fixed positive lower bound.  The
scalar version of \eqref{eq:damping-integrated} yields
\[
 \int_{\text{dwell }j}|w-d_{l,i}|\,dt
 \leq C+C L_jj^{-2}.
\]
The two ramps and the final rest have total duration \(3T_j\),
and their contribution to \(\int w\) is \(O(T_j)\).  Since the
solution starts at \(1\) and ends at \((D_{l,j})_{ii}\),
integration of \(w=(\log y)'\) gives the first equality in
\eqref{eq:diagonal-growth}.  The second follows from
\(T_j/L_j=j^{-1}\to0\) and \(L_j\to\infty\).  All constants
are uniform over the fixed finite list of coordinates.
\end{proof}

\subsection{Tuning every complete period exactly}
\label{subsec:tuning}

The preceding identity puts the actual full-period map in the
same local coordinate problem as the ideal word.  Its remaining
error tends to zero with one parameter derivative, uniformly
over the past.  This allows the target cycle to be attained
successively in every period.

\begin{proposition}[Exact full-period transmission]
\label{prop:transmission}
For finite spectral data as specified at the end of Section~\ref{sec:counting},
choose the word and compact pulse family above. There are a starting
index \(j_0\) and parameters \(h_j\in K\), for every \(j\geq j_0\),
such that the center-regular value transmissions of the resulting metric
satisfy
\begin{equation}
 \mathcal T_{l,j}=s_{l,j}D_{l,j}\Pi_l,\qquad
 s_{l,j}>0,\qquad |\log s_{l,j}|\leq C T_j
 \quad(2\leq l\leq L,\ j\geq j_0).
 \label{eq:exact-transmission}
\end{equation}
The positive diagonal matrices \(D_{l,j}\) propagate the reference
Cauchy data \((I,\theta_l I)\) over the entire leg following the pulse;
they satisfy \eqref{eq:diagonal-growth}. The regular value equation has the
bounds \eqref{eq:riccati-barriers}. All estimates used to choose \(j_0\)
are uniform over preceding admitted controls for this fixed finite family.
\end{proposition}

\begin{proof}
First fix an admitted past and vary the current parameter \(h\).
By \eqref{eq:full-normalization},
\[
 e^{-\varphi_{l,j}}D_{l,j}^{-1}\mathcal T_{l,j}(h)
 =\left[I+K_{l,j}\bigl(P_{l,e}(h)-\theta_lI\bigr)\right]
       e^{-\varphi_{l,j}}\mathcal B_{l,j}(h).
\]
The bound \eqref{eq:K-bound} and Lemma~\ref{lem:pulse} imply
\begin{equation}
 e^{-\varphi_{l,j}}D_{l,j}^{-1}\mathcal T_{l,j}(h)
 =H_l(1,h)+O_{C^1}(\eta_j).
 \label{eq:full-period-limit}
\end{equation}
The estimate is uniform in the preceding history.  Each
\(\mathcal T_{l,j}\) is a value fundamental matrix of
\(Y'=P_lY\), so it has positive determinant; so does \(D_{l,j}\).
Determinant normalization is therefore defined for their
quotient.  The matrices \(H_l(1,h)\), \(h\in K\), form a compact
subset of the positive-determinant matrices.  On a neighborhood
of this compact set, \(\mathcal N_l\) is smooth with bounded
derivatives.  Normalization is unchanged by the positive scalar
\(e^{-\varphi_{l,j}}\), so \eqref{eq:full-period-limit} gives
\begin{equation}
 \Psi_j(h):=
 \bigl(\mathcal N_l(D_{l,j}^{-1}\mathcal T_{l,j}(h))\bigr)_{l=2}^{L}
 =\Phi(h)+O_{C^1}(\eta_j)
 \quad\text{in }\mathcal G.
 \label{eq:group-limit}
\end{equation}
The assertion is first an estimate for the ambient matrix
entries and their derivatives, and hence also in any fixed
smooth chart containing the compact image under consideration.

We spell out the uniform inversion.  Use Euclidean norms on the
parameter and coordinate spaces, and the associated operator norms
for their derivatives.  Choose a coordinate chart
\(\zeta\) near \(\Pi\) in \(\mathcal G\) with \(\zeta(\Pi)=0\),
and let \(f=\zeta\circ\Phi\) and \(A_0=Df(0)\).  By
\eqref{eq:ideal-word}, \(A_0\) is invertible.  Choose a closed
ball \(\overline B_r\) centered at \(0\), contained in the
interior of \(K\), whose \(\Phi\)-image lies compactly inside
the chart and for which
\[
 \sup_{h\in\overline B_r}
 \|I-A_0^{-1}Df(h)\|\leq\frac14.
\]
For every sufficiently large \(j\), uniformly in the admitted
past, \(f_j=\zeta\circ\Psi_j\) is defined on this ball and
\[
 \sup_{h\in\overline B_r}
 \|A_0^{-1}(Df_j(h)-Df(h))\|\leq\frac14,\qquad
 \|A_0^{-1}f_j(0)\|\leq\frac r2.
\]
This follows from \eqref{eq:group-limit} and \(\eta_j\to0\).
The map
\[
 \mathcal C_j(h)=h-A_0^{-1}f_j(h)
\]
has derivative of norm at most \(1/2\) on the ball.  Also
\[
 \|\mathcal C_j(h)\|
 \leq\|\mathcal C_j(0)\|+\tfrac12\|h\|
 \leq r
 \qquad(h\in\overline B_r).
\]
It is therefore a contraction of the closed ball into itself.
Its iterates converge to a fixed point \(h_j\), for which
\(f_j(h_j)=0\), or \(\Psi_j(h_j)=\Pi\).

All the preceding estimates were uniform over the past.
Choose \(j_0\) once, large enough for them and for the geometric
input, after the fixed word and parameter ball have been chosen.
The history before \(j_0\) is round.  Choose \(h_{j_0}\) by the
contraction above, and repeat after each finite chosen history.
This inductively defines an infinite admitted sequence and
hence a single smooth metric of Proposition~\ref{prop:metric}.
Undoing determinant normalization at each fixed point gives
\[
 D_{l,j}^{-1}\mathcal T_{l,j}=s_{l,j}\Pi_l,\qquad
 s_{l,j}=\det(D_{l,j}^{-1}\mathcal T_{l,j})^{1/N_l}>0.
\]
By \eqref{eq:full-period-limit}, the determinant of
\(e^{-\varphi_{l,j}}D_{l,j}^{-1}\mathcal T_{l,j}\) stays
bounded above and away from zero on the fixed compact family.
Thus
\[
 \log s_{l,j}=\varphi_{l,j}+O(1)=O(T_j),
\]
because \(v_l\) is uniformly bounded.  This proves
\eqref{eq:exact-transmission}.  The equality is for the complete
period; no residual change of direction is passed to a later
period.

\end{proof}

\section{Global growth and cone geometry}
\label{sec:conclusion}

All parameters, the finite degree range, and the exact sequence of controls
are now fixed. We first prove that the selected center-regular solutions
satisfy the pointwise bound in Theorem~\ref{thm:main}. We then examine the
cone limits, which explain why additional hypotheses in the comparison
literature do not apply.

\subsection{Every-point growth and the strict count}
\begin{proof}[Proof of Theorem~\ref{thm:main}]
We convert the exact cycles from Proposition~\ref{prop:transmission} into polynomial growth.
For each \(l\) with \(M_l>0\) and each \(1\leq i\leq M_l\),
choose the center-regular solution with value vector
\(Y(t_{j_0})=e_{l,i}\).  Lemma~\ref{lem:regular} permits this
choice.  Near the origin it is \(r^l\) times an element of
\(V_l\), hence a homogeneous harmonic polynomial in Euclidean
coordinates.  It is smooth and harmonic there and, by
\eqref{eq:harmonic-ode}, on the rest of the manifold.

Let \(\sigma_l\) be the cyclic permutation of
\(\{1,\ldots,M_l\}\) underlying \(\Pi_l\), and define
\(i_{j_0}=i\), \(i_{j+1}=\sigma_l(i_j)\).  Applying
\eqref{eq:exact-transmission} sends \(e_{l,i_j}\) to a signed
multiple of \(e_{l,i_{j+1}}\), with absolute multiplier
\(s_{l,j}(D_{l,j})_{i_{j+1},i_{j+1}}\).  Hence, with
\[
 \varepsilon_{l,i,j}
   =L_j^{-1}\log(D_{l,j})_{ii}-d_{l,i},
 \qquad \max_{l,i}|\varepsilon_{l,i,j}|\longrightarrow0
\]
by \eqref{eq:diagonal-growth}, one has
\begin{equation}
 \log\|Y(t_J)\|=
 \sum_{j=j_0}^{J-1}
 \left(\log s_{l,j}
       +j^7\bigl(d_{l,i_{j+1}}+\varepsilon_{l,i_{j+1},j}\bigr)\right).
 \label{eq:cycle-sum}
\end{equation}
Signs in \(\Pi_l\) do not affect this norm.

We verify the weighted average in \eqref{eq:cycle-sum}.  Group
the sum of \(j^7d_{l,i_{j+1}}\) into complete cycles of length
\(M_l\).  Within a cycle starting at \(j\),
\((j+r)^7-j^7=O(j^6)\) for \(0\leq r<M_l\), with a constant
depending on this fixed length.  Since the sum of
\(d_{l,i_{j+r+1}}-\overline d_l\) in a complete cycle is zero,
that cycle contributes only \(O(j^6)\) to the difference from
the mean-weighted sum.  Summing these errors and bounding the
at most \(M_l-1\) final terms by \(O(J^7)\) gives
\[
 \sum_{j=j_0}^{J-1}j^7d_{l,i_{j+1}}
 =\overline d_l\sum_{j=j_0}^{J-1}j^7+O(J^7).
\]
Also \(\sum_{j<J}|\log s_{l,j}|=O(\sum_{j<J}j^6)=O(J^7)\).
The \(\varepsilon\)-terms sum to \(o(J^8)\): for any positive
number, choose a fixed late index after which their absolute
values are below that number, and use
\(\sum_{j<J}j^7=O(J^8)\).  The finitely many earlier terms
vanish after division by \(J^8\).  Finally,
\[
 \sum_{j=j_0}^{J-1}j^7=\frac{J^8}{8}+O(J^7),\qquad
 t_J=\frac{J^8}{8}+O(J^7)
\]
by \eqref{eq:periods}.  Therefore
\begin{equation}
 \lim_{J\to\infty}\frac{\log\|Y(t_J)\|}{t_J}
 =\overline d_l<k.
 \label{eq:endpoint-rate}
\end{equation}

The estimate extends to every radius.  From \(Y'=P_lY\) and
\eqref{eq:riccati-barriers},
\[
 \|Y(t)\|\leq
 \exp\bigl(C_P(t-t_j)\bigr)\|Y(t_j)\|
 \qquad(t_j\leq t\leq t_{j+1}).
\]
The ratio \((t_{j+1}-t_j)/t_j\) tends to zero by
\eqref{eq:time-asymptotic}.  Thus \eqref{eq:endpoint-rate} gives
\[
 \limsup_{t\to\infty}\frac{\log\|Y(t)\|}{t}
 \leq\overline d_l.
\]
Choose \(\alpha_l\) with
\(\overline d_l<\alpha_l<k\).  For all sufficiently large
\(r=e^t\), \(\|Y(\log r)\|\leq r^{\alpha_l}\).
The fixed finite-dimensional space \(V_l\) has
\[
 \left|\sum_iY_i e_{l,i}(\omega)\right|\leq C_l\|Y\|
 \quad\text{for every }\omega\in S^m,
\]
for example by Cauchy--Schwarz and boundedness of the finitely
many smooth basis functions.  Since \(d_g(0,x)=r(x)\), this
gives \(|u(x)|\leq C_u(1+d_g(0,x))^k\) outside a compact set.
Smoothness bounds \(u\) on that compact set and supplies the
same inequality everywhere.

These functions are linearly independent: their restrictions
to the sphere \(t=t_{j_0}\) are the selected vectors of
mutually orthogonal round spaces \(V_l\).  It follows that
\[
 \dim_{\mathbb R}\mathcal H_k(\mathbb R^{m+1},g)
 \geq\sum_{l=2}^{L}M_l>\sum_{l=0}^{k}N_l.
\]
The Euclidean identification was proved at the end of
Section~\ref{sec:counting}, so this is the required strict comparison.

Finally, the distance and volume statements in Proposition~\ref{prop:metric} give
\[
 \operatorname{vol}_g B_g(0,R)
 =\operatorname{vol}_{g_0}(S^m)
       \int_0^R r^m a(\log r)^m\,dr.
\]
Because \(a(t)\to a_*\), changing variables \(r=Rz\) and
using dominated convergence shows that this quantity divided
by \(\omega_{m+1}R^{m+1}\) tends to \(a_*^m\); here
\(\operatorname{vol}_{g_0}(S^m)=(m+1)\omega_{m+1}\).
This completes the proof.
\end{proof}

\subsection{Cone limits, conformal curvature, and cone degrees}
\begin{proposition}[Cone geometry and the union of degrees]
\label{prop:cone-consequences}
Let \(g\) be the metric constructed in the proof of Theorem~\ref{thm:main},
with link parameters \(a_*,q_*\) and comparison integer \(k\).
Every cone with link \(G(a_*,q,J_0)\), \(1\leq q\leq q_*\), occurs
as a pointed tangent cone at infinity. In particular the tangent cone is
not unique. The metric \(g\) is not locally conformally flat. If
\(\mathcal D(M)\) denotes the union of nonnegative homogeneous harmonic
degrees over all its tangent cones, then \(\mathcal D(M)\) contains a
half-line and a neighborhood of \(k\).
\end{proposition}

\begin{proof}
 Centering a fixed
logarithmic window in the round rests or nonround dwells and letting its
center tend to infinity gives smooth annular limits with links
$G(a_*,1,J_0)$ and $G(a_*,q_*,J_0)$, respectively. They are nonisometric:
the first is Einstein, whereas the second has distinct horizontal and
vertical Ricci eigenvalues. Every intermediate ramp link also occurs as a
limit, because the ramp durations diverge and their derivatives tend to
zero. These annular limits extend to pointed cone limits including the
vertex: the inner radial ball of radius $\delta$ has diameter at most
$2\delta$ in every rescaling and in the limit. Comparing paths on the
complementary annuli, then letting $\delta\downarrow0$, gives convergence
of distances on each bounded ball. All these links have the same volume.

The additional local conformal flatness hypothesis in~\cite{CaiLai}
also fails.  Put \(h_*=G(a_*,q_*,J_0)\).
The nonround cone metric \(dr^2+r^2h_*\) is conformal, under
\(t=\log r\), to the product \(dt^2+h_*\).
For this product the Weyl tensor has mixed component
\[
 W_{0a0b}
 =-\frac{1}{m-1}
   \left(\Ric_{h_*}-\frac{\operatorname{Scal}_{h_*}}{m}h_*\right)_{ab}.
\]
Indeed the product has zero mixed curvature and radial Ricci tensor;
substitution in the trace decomposition of curvature gives this formula.
The component is nonzero because the link has distinct horizontal and
vertical Ricci eigenvalues.  Vanishing of the Weyl tensor is conformally
invariant and is preserved by smooth annular limits.  If the constructed
metric were locally conformally flat, its nonround cone limit would have
zero Weyl tensor, contradicting the formula.

For maximal Hopf charge in degree $l$, its
growth root $\alpha_l(q)$ is the nonnegative solution of
\[
 \alpha_l(q)(\alpha_l(q)+m-1)
 =a_*^{-2}q^{1/m}
   \bigl[l(l+m-1)+(q^{-1}-1)l^2\bigr].
\]
This root varies continuously with $q$. At $q=1$ its asymptotic slope is
$a_*^{-1}$; at $q=q_*$ it is
$a_*^{-1}q_*^{-(m-1)/(2m)}$, which is strictly smaller. More precisely, the endpoint roots are
\[
 \alpha_l(1)=a_*^{-1}l+O(1),\qquad
 \alpha_l(q_*)=a_*^{-1}q_*^{-(m-1)/(2m)}l+O(1).
\]
The difference of their slopes is positive, so for every sufficiently
large \(l\), \(\alpha_{l+1}(q_*)<\alpha_l(1)\) and
\(\alpha_l(q_*)<\alpha_l(1)\). The continuous image of the
\(q\)-interval contains the entire interval between these endpoints.
These intervals overlap and have unbounded upper endpoints, so their
union contains a half-line. In particular, equal cone-link
volumes do not provide a discrete union spectrum.

The chosen integer $k$ itself lies in this union. Indeed,
\eqref{eq:selection} forces some selected degree $l>k$. Its smallest dwell
root is below $k$, while its round root $\theta_l$ is greater than $l$
because $a_*<1$. Continuity supplies an intermediate cone root equal to
$k$; the strict endpoint inequalities in fact put a neighborhood of $k$
in $\mathcal D(M)$. Xu's nonunique-cone three-circles
theorem~\cite[Theorem~3.4]{Xu} requires the comparison exponent to lie
outside $\mathcal D(M)$, so it cannot be applied at $k$ or at any
sufficiently nearby exponent. The unique-cone spectral bounds in~\cite{Huang} also
have a hypothesis that fails. A global logarithmic growth average, as
in~\eqref{eq:endpoint-rate}, need not equal the local homogeneous degree
seen after normalization on each separate cone subsequence.
\end{proof}

\appendix
\section{An exact finite spectral selection}
\label{sec:finite-selection}

The analytic argument in Lemma~\ref{lem:spectral-surplus} supplies all
spectral data needed in the proof of Theorem~\ref{thm:main}.  A supplementary
integer calculation gives a concrete selection in ambient dimension
\(16\), at the integer cutoff \(50000\).  This calculation checks the
finite spectral inequalities in~\eqref{eq:selection}; it does not construct
the finite control word or the metric numerically.  The analytic existence
argument remains independent of it.

The parameters of the calculation are
\begin{gather*}
 m=15,\qquad s=8,\qquad k=50000,\qquad q_*=\frac{101}{100},\\
 c=\frac{147}{1000},\qquad \alpha_*^2=\frac{99853}{100000},
 \qquad a_*^2=q_*^{1/15}\alpha_*^2.
\end{gather*}
The endpoint curvature inequality is strict because
\[
 1-\frac{2(q_*-1)}{m-1}-\alpha_*^2=\frac{29}{700000}>0.
\]
The integer inequality
\(101\cdot99853^{15}<100\cdot100000^{15}\)
shows that \(q_*(\alpha_*^2)^{15}<1\), and hence \(a_*<1\).
The decrease of the curvature threshold in the proof of
Lemma~\ref{lem:spectral-surplus} gives the margin for every
\(q\in[1,q_*]\).
The radial scale also meets its required smallness condition:
\[
 -\log a_*\le-\tfrac12\log\alpha_*^2
 <\frac{1-\alpha_*^2}{2\alpha_*^2}<\frac1{1000}.
\]
The cutoff integral in Proposition~\ref{prop:metric} is at least
\(\int_4^\infty(1+t)^{-5/4}\,dt=4\cdot5^{-1/4}>2\).
It therefore gives \(\mu<1/2000\), in particular \(b\ge-1/4\).
These are parameter checks; the control and transmission arguments are
still the analytic arguments of Sections~\ref{sec:control}--\ref{sec:conclusion}.

Here is the arithmetic bound underlying the selection.  Write
\(h=m-1\), \(Q=10^8\), and an angular eigenvalue as
\(\lambda=\mathrm{num}/\mathrm{den}>0\), with integer numerator and
denominator.  If \(\operatorname{isqrt}(a)=\lfloor\sqrt a\rfloor\) for a
nonnegative integer \(a\), define
\[
 S=\operatorname{isqrt}\!\left(
    \left\lfloor\frac{Q^2(h^2\mathrm{den}+4\mathrm{num})}
                         {\mathrm{den}}\right\rfloor\right)+1.
\]
Then \(S^2\mathrm{den}>Q^2(h^2\mathrm{den}+4\mathrm{num})\), so the
positive root of \(d(d+h)=\lambda\) satisfies the strict rational bound
\[
 d<\frac{S-hQ}{2Q}.
\]
For the charge indexed by \(b\), the exact eigenvalue data are
\[
 \mathrm{num}=100000\bigl(101l(l+14)-(2b-l)^2\bigr),
 \qquad \mathrm{den}=101\cdot99853.
\]
The multiplicities are the integers in~\eqref{eq:hopf-multiplicity}.
The two equal contributions indexed by \(b\) and \(l-b\) are combined,
except for the middle term when they coincide.  Increasing \(b\) from
\(0\) to \(\lfloor l/2\rfloor\) puts the eigenvalues, and therefore the
upper bounds for the exponents, in increasing order.

The supplied program \texttt{verification/counting-check.py}
uses these integer multiplicities and upper bounds.  It accepts an initial
group only when its upper-bounded exponent sum is strictly less than
\(k\) times its size.  It may take part of a repeated eigenvalue, which
corresponds to choosing that many vectors in its real eigenspace.  Low
degrees are selected in full whenever the exact inequality
\(l(l+14)<\alpha_*^2k(k+14)\) certifies all their roots to be below \(k\).
In odd degree the zero Hopf charge is absent, so the full-degree test
is a sufficient criterion rather than an exact characterization of the
largest root. The individual search ends when the smallest root is at least \(k\).
Thus every counted group has a true mean strictly below \(k\).

The resulting counts, using integer arithmetic, are
\begin{align*}
 \text{selected dimension}
   &=46785605044474340387088811257657817010714687366486217934111,\\
 h_{50000}(\mathbb R^{16})
   &=46779709349146362239266126538609505511855492134986686636251,\\
 \text{strict surplus}
   &=5895695327978147822684719048311498859195231499531297860.
\end{align*}
The selected dimension omits degrees \(0\) and \(1\), as required in
\eqref{eq:selection}.  Only the displayed strict integer comparison is
used by the certificate; a decimal ratio is unnecessary.

These parameters and selections satisfy all the finite-data hypotheses at
the end of Section~\ref{sec:counting}. Applying
Propositions~\ref{prop:metric} and~\ref{prop:transmission}, followed by the
growth argument in Section~\ref{sec:conclusion}, therefore gives the metric
of Theorem~\ref{thm:main} with \(n=16\) and \(k=50000\).

\clearpage
\begin{thebibliography}{99}

\bibitem{BourguignonKarcher}
J.-P.~Bourguignon and H.~Karcher,
\emph{Curvature operators: pinching estimates and geometric examples},
Ann. Sci. \'Ecole Norm. Sup. (4) \textbf{11} (1978), no.~1, 71--92.
\href{https://doi.org/10.24033/asens.1340}{doi:10.24033/asens.1340}.

\bibitem{CaiLai}
X.~Cai and M.~Lai,
\emph{Dimension of polynomial growth harmonic functions on locally
conformally flat manifolds with nonnegative Ricci curvature},
preprint, \href{https://arxiv.org/abs/2608.04553v2}{arXiv:2608.04553v2}, 2026.

\bibitem{CM97}
T.~H.~Colding and W.~P.~Minicozzi~II,
\emph{Harmonic functions on manifolds},
Ann. of Math. (2) \textbf{146} (1997), no.~3, 725--747.
\href{https://doi.org/10.2307/2952459}{doi:10.2307/2952459}.

\bibitem{CM98}
T.~H.~Colding and W.~P.~Minicozzi~II,
\emph{Weyl type bounds for harmonic functions},
Invent. Math. \textbf{131} (1998), no.~2, 257--298.
\href{https://doi.org/10.1007/s002220050204}{doi:10.1007/s002220050204}.

\bibitem{ColdingNaber}
T.~H.~Colding and A.~Naber,
\emph{Characterization of tangent cones of noncollapsed limits with lower
Ricci bounds and applications},
Geom. Funct. Anal. \textbf{23} (2013), no.~1, 134--148.
\href{https://doi.org/10.1007/s00039-012-0202-7}{doi:10.1007/s00039-012-0202-7}.
Author preprint: \href{https://arxiv.org/abs/1108.3244v3}{arXiv:1108.3244v3}
(cited lemma numbering follows this version).

\bibitem{Donnelly}
H.~Donnelly,
\emph{Harmonic functions on manifolds of nonnegative Ricci curvature},
Int. Math. Res. Not. \textbf{2001} (2001), no.~8, 429--434.
\href{https://doi.org/10.1155/S1073792801000216}{doi:10.1155/S1073792801000216}.

\bibitem{FigueroaUngureanu2016}
J.~Figueroa-O'Farrill and M.~Ungureanu,
\emph{Homogeneous M2 duals},
J. High Energy Phys. \textbf{2016} (2016), no.~1, article~150.
\href{https://doi.org/10.1007/JHEP01(2016)150}{doi:10.1007/JHEP01(2016)150}.

\bibitem{Huang}
X.-T.~Huang,
\emph{Harmonic functions with polynomial growth on manifolds with
nonnegative Ricci curvature},
Calc. Var. Partial Differential Equations \textbf{62} (2023),
article~111.
\href{https://doi.org/10.1007/s00526-023-02456-z}{doi:10.1007/s00526-023-02456-z}.
Author preprint: \href{https://arxiv.org/abs/2109.07534v1}{arXiv:2109.07534v1}
(cited theorem numbering follows this version).

\bibitem{JurdjevicSussmann}
V.~Jurdjevic and H.~J.~Sussmann,
\emph{Control systems on Lie groups},
J. Differential Equations \textbf{12} (1972), no.~2, 313--329.
\href{https://doi.org/10.1016/0022-0396(72)90035-6}{doi:10.1016/0022-0396(72)90035-6}.

\bibitem{LauretMiatelloRossetti}
E.~A.~Lauret, R.~J.~Miatello and J.~P.~Rossetti,
\emph{Spectra of lens spaces from $1$-norm spectra of congruence lattices},
Int. Math. Res. Not. \textbf{2016} (2016), no.~4, 1054--1089.
\href{https://doi.org/10.1093/imrn/rnv159}{doi:10.1093/imrn/rnv159}.
Author preprint: \href{https://arxiv.org/abs/1311.7167v4}{arXiv:1311.7167v4}
(cited lemma numbering follows this version).

\bibitem{LawsonLim}
J.~Lawson and Y.~Lim,
\emph{A Birkhoff contraction formula with applications to Riccati equations},
SIAM J. Control Optim. \textbf{46} (2007), no.~3, 930--951.
\href{https://doi.org/10.1137/050637637}{doi:10.1137/050637637}.

\bibitem{Lee1992}
C.-I.~C.~Lee,
\emph{On Laplace continued fraction for the normal integral},
Ann. Inst. Statist. Math. \textbf{44} (1992), no.~1, 107--120.
\href{https://www.ism.ac.jp/editsec/aism/pdf/044_1_0107.pdf}{Published text}.

\bibitem{LiTam}
P.~Li and L.-F.~Tam,
\emph{Linear growth harmonic functions on a complete manifold},
J. Differential Geom. \textbf{29} (1989), no.~2, 421--425.
\href{https://doi.org/10.4310/jdg/1214442883}{doi:10.4310/jdg/1214442883}.

\bibitem{Perelman}
G.~Perelman,
\emph{A complete Riemannian manifold of positive Ricci curvature with
Euclidean volume growth and nonunique asymptotic cone},
in \emph{Comparison Geometry},
Math. Sci. Res. Inst. Publ. \textbf{30}, Cambridge University Press,
1997, pp.~165--166.
\href{https://library.slmath.org/books/Book30/files/perex.pdf}{Published text}.

\bibitem{Tanno}
S.~Tanno,
\emph{The first eigenvalue of the Laplacian on spheres},
Tohoku Math. J. (2) \textbf{31} (1979), no.~2, 179--185.
\href{https://doi.org/10.2748/tmj/1178229837}{doi:10.2748/tmj/1178229837}.

\bibitem{Xu}
G.~Xu,
\emph{Three circles theorems for harmonic functions},
Math. Ann. \textbf{366} (2016), nos.~3--4, 1281--1317.
\href{https://doi.org/10.1007/s00208-016-1366-5}{doi:10.1007/s00208-016-1366-5}.
Author preprint: \href{https://arxiv.org/abs/1601.02066v1}{arXiv:1601.02066v1}
(cited theorem numbering follows this version).

\bibitem{Yau}
S.-T.~Yau,
\emph{Open problems in geometry},
in R.~Schoen and S.-T.~Yau,
\emph{Lectures on Differential Geometry},
International Press, 1994, pp.~365--404, Problem~48.

\end{thebibliography}

\end{document}
